\documentclass[11pt]{article}

% ---- theme variables (passed via pandoc -V) ----
\newcommand{\ThemeName}{Transformers \& LLMs}
\newcommand{\ThemeKicker}{AI/ML Track}
\newcommand{\ThemeNumber}{04}

\usepackage{interviewstyle}
\setthemecolor{7c3aed}{a78bfa}

% pandoc-required plumbing
\providecommand{\tightlist}{\setlength{\itemsep}{1pt}\setlength{\parskip}{0pt}}
\setlength{\emergencystretch}{3em}
\providecommand{\pandocbounded}[1]{#1}

% syntax-highlighting macros emitted by pandoc (skylighting)
\usepackage{color}
\usepackage{fancyvrb}
\newcommand{\VerbBar}{|}
\newcommand{\VERB}{\Verb[commandchars=\\\{\}]}
\DefineVerbatimEnvironment{Highlighting}{Verbatim}{commandchars=\\\{\}}
% Add ',fontsize=\small' for more characters per line
\usepackage{framed}
\definecolor{shadecolor}{RGB}{35,38,41}
\newenvironment{Shaded}{\begin{snugshade}}{\end{snugshade}}
\newcommand{\AlertTok}[1]{\textcolor[rgb]{0.58,0.85,0.30}{\textbf{\colorbox[rgb]{0.30,0.12,0.14}{#1}}}}
\newcommand{\AnnotationTok}[1]{\textcolor[rgb]{0.25,0.50,0.35}{#1}}
\newcommand{\AttributeTok}[1]{\textcolor[rgb]{0.16,0.50,0.73}{#1}}
\newcommand{\BaseNTok}[1]{\textcolor[rgb]{0.96,0.45,0.00}{#1}}
\newcommand{\BuiltInTok}[1]{\textcolor[rgb]{0.50,0.55,0.55}{#1}}
\newcommand{\CharTok}[1]{\textcolor[rgb]{0.24,0.68,0.91}{#1}}
\newcommand{\CommentTok}[1]{\textcolor[rgb]{0.48,0.49,0.49}{#1}}
\newcommand{\CommentVarTok}[1]{\textcolor[rgb]{0.50,0.55,0.55}{#1}}
\newcommand{\ConstantTok}[1]{\textcolor[rgb]{0.15,0.68,0.68}{\textbf{#1}}}
\newcommand{\ControlFlowTok}[1]{\textcolor[rgb]{0.99,0.74,0.29}{\textbf{#1}}}
\newcommand{\DataTypeTok}[1]{\textcolor[rgb]{0.16,0.50,0.73}{#1}}
\newcommand{\DecValTok}[1]{\textcolor[rgb]{0.96,0.45,0.00}{#1}}
\newcommand{\DocumentationTok}[1]{\textcolor[rgb]{0.64,0.20,0.25}{#1}}
\newcommand{\ErrorTok}[1]{\textcolor[rgb]{0.85,0.27,0.33}{\underline{#1}}}
\newcommand{\ExtensionTok}[1]{\textcolor[rgb]{0.00,0.60,1.00}{\textbf{#1}}}
\newcommand{\FloatTok}[1]{\textcolor[rgb]{0.96,0.45,0.00}{#1}}
\newcommand{\FunctionTok}[1]{\textcolor[rgb]{0.56,0.27,0.68}{#1}}
\newcommand{\ImportTok}[1]{\textcolor[rgb]{0.15,0.68,0.38}{#1}}
\newcommand{\InformationTok}[1]{\textcolor[rgb]{0.77,0.36,0.00}{#1}}
\newcommand{\KeywordTok}[1]{\textcolor[rgb]{0.81,0.81,0.76}{\textbf{#1}}}
\newcommand{\NormalTok}[1]{\textcolor[rgb]{0.81,0.81,0.76}{#1}}
\newcommand{\OperatorTok}[1]{\textcolor[rgb]{0.81,0.81,0.76}{#1}}
\newcommand{\OtherTok}[1]{\textcolor[rgb]{0.15,0.68,0.38}{#1}}
\newcommand{\PreprocessorTok}[1]{\textcolor[rgb]{0.15,0.68,0.38}{#1}}
\newcommand{\RegionMarkerTok}[1]{\textcolor[rgb]{0.16,0.50,0.73}{\colorbox[rgb]{0.08,0.19,0.26}{#1}}}
\newcommand{\SpecialCharTok}[1]{\textcolor[rgb]{0.24,0.68,0.91}{#1}}
\newcommand{\SpecialStringTok}[1]{\textcolor[rgb]{0.85,0.27,0.33}{#1}}
\newcommand{\StringTok}[1]{\textcolor[rgb]{0.96,0.31,0.31}{#1}}
\newcommand{\VariableTok}[1]{\textcolor[rgb]{0.15,0.68,0.68}{#1}}
\newcommand{\VerbatimStringTok}[1]{\textcolor[rgb]{0.85,0.27,0.33}{#1}}
\newcommand{\WarningTok}[1]{\textcolor[rgb]{0.85,0.27,0.33}{#1}}

\begin{document}

\makecoverpage

\thispagestyle{empty}
{\hypersetup{hidelinks}\tableofcontents}
\clearpage

\begin{quote}
The Transformer replaced recurrence with attention, and that single architectural decision unlocked the scaling laws that produced GPT-4, Claude, Gemini, and every frontier AI system in production today.
\end{quote}

\section{Overview --- What it is}\label{overview--what-it-is}

A \textbf{Transformer} is a neural network architecture built entirely on attention mechanisms --- no recurrence, no convolution. Introduced by Vaswani et al. in 2017 ("Attention Is All You Need"), it processes entire sequences in parallel, allowing massive parallelism on modern GPUs/TPUs.

A \textbf{Large Language Model (LLM)} is a Transformer trained at scale (billions of parameters, trillions of tokens) to model the probability distribution of text. The emergent result is a system that can reason, code, translate, summarize, answer questions, and engage in open-ended dialogue.

\subsection{What makes them special}\label{what-makes-them-special}

\begin{itemize}
\tightlist
\item
  \textbf{Attention is global}: every token can directly attend to every other token in the context
\item
  \textbf{Parallelizable training}: unlike RNNs, all positions are processed simultaneously
\item
  \textbf{Scale works}: more parameters + more data + more compute = reliably better models (scaling laws)
\item
  \textbf{Transfer learning}: pretrain once on internet-scale data, fine-tune cheaply for any task
\end{itemize}

\subsection{Key numbers to know}\label{key-numbers-to-know}

\begin{longtable}[]{@{}
  >{\raggedright\arraybackslash}p{(\columnwidth - 6\tabcolsep) * \real{0.3455}}
  >{\raggedright\arraybackslash}p{(\columnwidth - 6\tabcolsep) * \real{0.2337}}
  >{\raggedright\arraybackslash}p{(\columnwidth - 6\tabcolsep) * \real{0.3272}}
  >{\raggedright\arraybackslash}p{(\columnwidth - 6\tabcolsep) * \real{0.0935}}@{}}
\toprule\noalign{}
\rowcolor{acc!16}
\begin{minipage}[b]{\linewidth}\raggedright
Model
\end{minipage} & \begin{minipage}[b]{\linewidth}\raggedright
Parameters
\end{minipage} & \begin{minipage}[b]{\linewidth}\raggedright
Context Window
\end{minipage} & \begin{minipage}[b]{\linewidth}\raggedright
Year
\end{minipage} \\
\midrule\noalign{}
\endhead
\bottomrule\noalign{}
\endlastfoot
BERT-base & 110M & 512 tokens & 2018 \\
\rowcolor{zebra}
GPT-2 & 1.5B & 1024 tokens & 2019 \\
GPT-3 & 175B & 2048 tokens & 2020 \\
\rowcolor{zebra}
PaLM & 540B & 2048 tokens & 2022 \\
GPT-4 (est.) & \textasciitilde1.8T (MoE) & 128K tokens & 2023 \\
\rowcolor{zebra}
Llama 3.1 & 405B & 128K tokens & 2024 \\
Claude 3.5 Sonnet & Unknown & 200K tokens & 2024 \\
\rowcolor{zebra}
Gemini 1.5 Pro & Unknown & 1M tokens & 2024 \\
\end{longtable}

\section{\texorpdfstring{Why It Exists (the path from RNNs \(\rightarrow\) "Attention Is All You Need" \(\rightarrow\) GPT scale)}{Why It Exists (the path from RNNs \textbackslash rightarrow "Attention Is All You Need" \textbackslash rightarrow GPT scale)}}\label{why-it-exists-the-path-from-rnns-rightarrow-attention-is-all-you-need-rightarrow-gpt-scale}

\subsection{The problem with RNNs and LSTMs}\label{the-problem-with-rnns-and-lstms}

Before Transformers, sequence modeling relied on \textbf{Recurrent Neural Networks (RNNs)} and their improved variant \textbf{LSTMs} (Long Short-Term Memory):

\setcodelang{}

\begin{verbatim}
RNN: h_t = tanh(W_h * h_{t-1} + W_x * x_t + b)

Input:  [w1] → [w2] → [w3] → [w4] → [w5]
           ↓      ↓      ↓      ↓      ↓
Hidden: [h1] → [h2] → [h3] → [h4] → [h5]
\end{verbatim}

\textbf{Fatal flaws:}

\begin{enumerate}
\def\labelenumi{\arabic{enumi}.}
\tightlist
\item
  \textbf{Sequential bottleneck}: h\_t depends on h\_\{t-1\} --- cannot parallelize across time steps
\item
  \textbf{Vanishing gradient}: gradients decay exponentially over long sequences; LSTM helped but didn\textquotesingle t solve it
\item
  \textbf{Fixed-size bottleneck}: the entire past is compressed into a single hidden vector h\_t
\item
  \textbf{Long-range forgetting}: "The cat, which sat on the mat, was..." \(\rightarrow\) by the time you reach "was", the model has weakened signal from "cat"
\end{enumerate}

\subsection{Seq2Seq + Bahdanau Attention (2014-2015) --- the precursor}\label{seq2seq--bahdanau-attention-2014-2015--the-precursor}

Bahdanau et al. added an \textbf{attention mechanism} to seq2seq models: the decoder, when generating token t, could look back at ALL encoder hidden states and form a weighted sum. This was a breakthrough --- but attention was still bolted onto recurrent models.

\setcodelang{Python}

\begin{Shaded}
\begin{Highlighting}[]
\NormalTok{Attention score: e\_\{ij\} }\OperatorTok{=}\NormalTok{ alignment(s\_\{i}\OperatorTok{{-}}\DecValTok{1}\NormalTok{\}, h\_j)}
\NormalTok{                α\_\{ij\} }\OperatorTok{=}\NormalTok{ softmax(e\_\{ij\})}
\NormalTok{Context vector: c\_i }\OperatorTok{=}\NormalTok{ Σ\_j α\_\{ij\} }\OperatorTok{*}\NormalTok{ h\_j}
\end{Highlighting}
\end{Shaded}

\subsection{"Attention Is All You Need" (2017) --- the revolution}\label{attention-is-all-you-need-2017--the-revolution}

Vaswani et al. asked: \textbf{what if we use ONLY attention and remove recurrence entirely?}

Key insight: if attention can let any position see any other position, we don\textquotesingle t need sequential processing at all. Just stack attention layers with position information baked in.

Result: the Transformer.

\textbf{Why this was transformative:}

\begin{itemize}
\tightlist
\item
  Training became fully parallelizable \(\rightarrow\) 10-100x faster
\item
  No gradient vanishing over sequence length
\item
  Any token directly attends to any other token (O(1) path length vs O(n) for RNN)
\item
  Scales cleanly with more compute
\end{itemize}

\subsection{The GPT arc (2018-present)}\label{the-gpt-arc-2018-present}

\setcodelang{}

\begin{verbatim}
2018: GPT-1    (117M params)  → pretraining + fine-tuning works
2018: BERT     (340M params)  → bidirectional pretraining dominates NLU
2019: GPT-2    (1.5B params)  → "too dangerous to release" — language model can write
2020: GPT-3    (175B params)  → few-shot learning, emergent abilities
2021: Codex    (12B params)   → GitHub Copilot launches
2022: ChatGPT  (GPT-3.5+RLHF) → 100M users in 60 days
2023: GPT-4    (~1.8T MoE?)   → multimodal, bar exam top 10%
2024: Claude/Gemini/Llama 3 → long context, open weights, SOTA
\end{verbatim}

\textbf{The key realization}: the pretraining objective (predict next token) is an \emph{extremely} rich self-supervised signal. Given enough data and compute, models develop internal representations of grammar, facts, reasoning, code, and more --- without any labeled data.

\section{Why FAANG \& AI Labs Care}\label{why-faang--ai-labs-care}

\subsection{OpenAI}\label{openai}

\begin{itemize}
\tightlist
\item
  Core product IS the LLM (GPT-4, o1, o3). Revenue: API + ChatGPT subscriptions
\item
  Research frontier: RLHF, Constitutional AI concepts, reasoning (o1/o3 chain-of-thought)
\item
  Interview focus: attention internals, training at scale, alignment
\end{itemize}

\subsection{Anthropic}\label{anthropic}

\begin{itemize}
\tightlist
\item
  Safety-focused frontier lab. Claude family (Haiku/Sonnet/Opus)
\item
  Pioneered Constitutional AI (RLHF with a "constitution" of principles instead of human raters)
\item
  Research: mechanistic interpretability, long context (200K), computer use
\item
  Interview focus: RLHF vs DPO, safety, interpretability, scaling
\end{itemize}

\subsection{Google DeepMind}\label{google-deepmind}

\begin{itemize}
\tightlist
\item
  Gemini family (Ultra/Pro/Flash/Nano), Bard/Gemini app
\item
  TPU infrastructure, PaLM/Chinchilla scaling laws, AlphaCode
\item
  Research: Chinchilla (optimal compute allocation), Flash Attention, speculative decoding
\item
  Interview focus: distributed training (TP/PP/DP), TPU design, efficient inference
\end{itemize}

\subsection{Meta AI}\label{meta-ai}

\begin{itemize}
\tightlist
\item
  Llama family (open weights --- Llama 2, 3, 3.1 up to 405B)
\item
  FAISS for vector search, PyTorch foundation, LLaMA \(\rightarrow\) LLaMA Guard
\item
  Research: LoRA (Microsoft/Meta), quantization, RoPE positional encodings
\item
  Interview focus: open source ecosystem, PEFT methods, inference optimization
\end{itemize}

\subsection{Microsoft}\label{microsoft}

\begin{itemize}
\tightlist
\item
  Heavy Azure OpenAI investment, GitHub Copilot, Bing Chat
\item
  Research: LoRA (Hu et al. 2021), DeepSpeed (ZeRO optimizer), GPTQ quantization
\item
  Interview focus: serving LLMs at scale, LoRA, product integration
\end{itemize}

\subsection{Amazon (AWS)}\label{amazon-aws}

\begin{itemize}
\tightlist
\item
  Bedrock platform (multi-model API), Titan models, Alexa LLM
\item
  Focus: enterprise deployment, cost optimization, retrieval-augmented generation (RAG)
\item
  Interview focus: inference economics, RAG pipelines, MLOps, serving latency
\end{itemize}

\textbf{Bottom line}: every major tech company\textquotesingle s AI strategy depends on Transformers. Understanding them deeply is table stakes for any AI/ML role at FAANG level.

\section{Core Concepts}\label{core-concepts}

\subsection{Tokenization}\label{tokenization}

Before a model sees text, it must be converted to discrete tokens (integers).

\subsubsection{Why not characters or words?}\label{why-not-characters-or-words}

\begin{longtable}[]{@{}
  >{\raggedright\arraybackslash}p{(\columnwidth - 4\tabcolsep) * \real{0.1458}}
  >{\raggedright\arraybackslash}p{(\columnwidth - 4\tabcolsep) * \real{0.2515}}
  >{\raggedright\arraybackslash}p{(\columnwidth - 4\tabcolsep) * \real{0.6028}}@{}}
\toprule\noalign{}
\rowcolor{acc!16}
\begin{minipage}[b]{\linewidth}\raggedright
Unit
\end{minipage} & \begin{minipage}[b]{\linewidth}\raggedright
Vocabulary size
\end{minipage} & \begin{minipage}[b]{\linewidth}\raggedright
Problem
\end{minipage} \\
\midrule\noalign{}
\endhead
\bottomrule\noalign{}
\endlastfoot
Characters & \textasciitilde256 & Sequences too long; no semantic meaning \\
\rowcolor{zebra}
Words & 300K+ & Rare words missing; morphology ignored \\
Subwords & 30K-100K & Sweet spot --- handles OOV, reasonable length \\
\end{longtable}

\subsubsection{Byte-Pair Encoding (BPE)}\label{byte-pair-encoding-bpe}

The dominant approach (GPT family, RoBERTa):

\begin{enumerate}
\def\labelenumi{\arabic{enumi}.}
\tightlist
\item
  Start: each character is a token
\item
  Count all consecutive pair frequencies
\item
  Merge the most frequent pair into a new token
\item
  Repeat until vocabulary size reached
\end{enumerate}

\setcodelang{}

\begin{verbatim}
Corpus: "low lower lowest"

Start:  l o w , l o w e r , l o w e s t
Step 1: "lo" is frequent → merge: lo w , lo w e r , lo w e s t
Step 2: "low" is frequent → merge: low , low e r , low e s t
Step 3: "lower" is frequent → merge: low , lower , low e s t
...
\end{verbatim}

\textbf{Result}: common words become single tokens; rare words split into meaningful pieces.
\texttt{"ChatGPT"} \(\rightarrow\) \texttt{{[}"Chat",\ "G",\ "PT"{]}} (approximately)

\subsubsection{WordPiece (BERT)}\label{wordpiece-bert}

Similar to BPE but maximizes likelihood instead of frequency. Uses \texttt{\#\#} prefix for continuation pieces: \texttt{"playing"} \(\rightarrow\) \texttt{{[}"play",\ "\#\#ing"{]}}

\subsubsection{SentencePiece (T5, Llama)}\label{sentencepiece-t5-llama}

Language-agnostic, treats input as raw bytes/unicode. No pre-tokenization required. Supports unigram language model or BPE.

\textbf{Key interview point}: tokenization affects everything --- cost (you pay per token), context window utilization, and model behavior on numbers/code. Numbers are often split digit-by-digit: \texttt{"12345"} \(\rightarrow\) \texttt{{[}"1",\ "2",\ "3",\ "4",\ "5"{]}} --- this is why LLMs struggle with arithmetic.

\subsection{Embeddings}\label{embeddings}

After tokenization, each token ID is mapped to a dense vector via a learned \textbf{embedding matrix} E \(\in\) \(\mathbb{R}\)\^{}\{V\(\times\)d\_model\}.

\setcodelang{}

\begin{verbatim}
Token: "cat" → ID: 5023
Embedding: E[5023] = [0.23, -0.14, 0.87, ..., 0.02]  # d_model dimensions
\end{verbatim}

\textbf{Properties:}

\begin{itemize}
\tightlist
\item
  Semantic similarity encoded in vector space
\item
  \texttt{king\ -\ man\ +\ woman\ ≈\ queen} (classic example)
\item
  These embeddings are learned during pretraining and encode rich linguistic knowledge
\item
  Input to the Transformer is: token embeddings + positional encodings
\end{itemize}

\subsection{Self-Attention --- The Core Mechanism}\label{self-attention--the-core-mechanism}

\visualfigure{../figures/aiml-04-tf/01-attention.pdf}{Each query token forms a weighted average over all tokens via softmax(QK$^{\mathsf{T}}$/$\surd$d). Rows sum to 1 — the model learns which context positions matter.}

\textbf{The fundamental idea}: for each token, compute a weighted sum of all token representations, where weights reflect relevance.

\subsubsection{Query, Key, Value (Q, K, V)}\label{query-key-value-q-k-v}

Every token produces three vectors via learned linear projections:

\setcodelang{Python}

\begin{Shaded}
\begin{Highlighting}[]
\NormalTok{Q }\OperatorTok{=}\NormalTok{ X W\_Q    (Query:  }\StringTok{"What am I looking for?"}\NormalTok{)}
\NormalTok{K }\OperatorTok{=}\NormalTok{ X W\_K    (Key:    }\StringTok{"What do I offer/label myself as?"}\NormalTok{)}
\NormalTok{V }\OperatorTok{=}\NormalTok{ X W\_V    (Value:  }\StringTok{"What is my actual content?"}\NormalTok{)}

\NormalTok{Where X ∈ ℝ}\OperatorTok{\^{}}\NormalTok{\{n×d\_model\}, W\_Q, W\_K, W\_V ∈ ℝ}\OperatorTok{\^{}}\NormalTok{\{d\_model×d\_k\}}
\end{Highlighting}
\end{Shaded}

\textbf{Intuition (from the newsroom analogy):}

\begin{itemize}
\tightlist
\item
  \textbf{Q (Query)}: the question an editor is asking ("who is the subject of this sentence?")
\item
  \textbf{K (Key)}: each word\textquotesingle s label/tag on the board ("I am a proper noun", "I am a verb")
\item
  \textbf{V (Value)}: each word\textquotesingle s actual content (the word\textquotesingle s semantic meaning to pass along)
\end{itemize}

\subsubsection{Scaled Dot-Product Attention}\label{scaled-dot-product-attention}

\setcodelang{}

\begin{verbatim}
                          QK^T
Attention(Q, K, V) = softmax(------) V
                          √d_k
\end{verbatim}

Step-by-step:

\setcodelang{Python}

\begin{Shaded}
\begin{Highlighting}[]
\NormalTok{Step }\DecValTok{1}\NormalTok{: Compute raw scores (compatibility between each query }\KeywordTok{and} \BuiltInTok{all}\NormalTok{ keys)}
\NormalTok{        scores }\OperatorTok{=}\NormalTok{ Q · K}\OperatorTok{\^{}}\NormalTok{T          shape: (n\_tokens × n\_tokens)}

\NormalTok{Step }\DecValTok{2}\NormalTok{: Scale to prevent softmax saturation}
\NormalTok{        scores }\OperatorTok{=}\NormalTok{ scores }\OperatorTok{/}\NormalTok{ sqrt(d\_k)}

\NormalTok{Step }\DecValTok{3}\NormalTok{: Apply softmax to get attention weights (probabilities)}
\NormalTok{        weights }\OperatorTok{=}\NormalTok{ softmax(scores)   shape: (n\_tokens × n\_tokens)}
        \CommentTok{\# Each row sums to 1.0}

\NormalTok{Step }\DecValTok{4}\NormalTok{: Weighted }\BuiltInTok{sum}\NormalTok{ of values}
\NormalTok{        output }\OperatorTok{=}\NormalTok{ weights · V        shape: (n\_tokens × d\_k)}
\end{Highlighting}
\end{Shaded}

\textbf{Why \(\surd\)d\_k scaling?} Without it, for large d\_k, dot products grow large in magnitude, pushing softmax into saturation regions where gradients vanish. Dividing by \(\surd\)d\_k keeps the variance of the dot product \textasciitilde1.

\setcodelang{Python}

\begin{Shaded}
\begin{Highlighting}[]
\ImportTok{import}\NormalTok{ numpy }\ImportTok{as}\NormalTok{ np}

\KeywordTok{def}\NormalTok{ scaled\_dot\_product\_attention(Q, K, V, mask}\OperatorTok{=}\VariableTok{None}\NormalTok{):}
    \CommentTok{"""}
\CommentTok{    Q: (batch, heads, seq\_len, d\_k)}
\CommentTok{    K: (batch, heads, seq\_len, d\_k)}
\CommentTok{    V: (batch, heads, seq\_len, d\_v)}
\CommentTok{    """}
\NormalTok{    d\_k }\OperatorTok{=}\NormalTok{ Q.shape[}\OperatorTok{{-}}\DecValTok{1}\NormalTok{]}
    
    \CommentTok{\# (batch, heads, seq\_len, seq\_len)}
\NormalTok{    scores }\OperatorTok{=}\NormalTok{ np.matmul(Q, K.transpose(}\DecValTok{0}\NormalTok{, }\DecValTok{1}\NormalTok{, }\DecValTok{3}\NormalTok{, }\DecValTok{2}\NormalTok{)) }\OperatorTok{/}\NormalTok{ np.sqrt(d\_k)}
    
    \ControlFlowTok{if}\NormalTok{ mask }\KeywordTok{is} \KeywordTok{not} \VariableTok{None}\NormalTok{:}
\NormalTok{        scores }\OperatorTok{=}\NormalTok{ scores }\OperatorTok{+}\NormalTok{ mask  }\CommentTok{\# mask = {-}inf for positions to ignore}
    
\NormalTok{    weights }\OperatorTok{=}\NormalTok{ np.exp(scores) }\OperatorTok{/}\NormalTok{ np.}\BuiltInTok{sum}\NormalTok{(np.exp(scores), axis}\OperatorTok{={-}}\DecValTok{1}\NormalTok{, keepdims}\OperatorTok{=}\VariableTok{True}\NormalTok{)}
    
    \CommentTok{\# (batch, heads, seq\_len, d\_v)}
\NormalTok{    output }\OperatorTok{=}\NormalTok{ np.matmul(weights, V)}
    \ControlFlowTok{return}\NormalTok{ output, weights}
\end{Highlighting}
\end{Shaded}

\subsubsection{Attention Complexity}\label{attention-complexity}

\begin{longtable}[]{@{}
  >{\raggedright\arraybackslash}p{(\columnwidth - 2\tabcolsep) * \real{0.1731}}
  >{\raggedright\arraybackslash}p{(\columnwidth - 2\tabcolsep) * \real{0.8269}}@{}}
\toprule\noalign{}
\rowcolor{acc!16}
\begin{minipage}[b]{\linewidth}\raggedright
Dimension
\end{minipage} & \begin{minipage}[b]{\linewidth}\raggedright
Cost
\end{minipage} \\
\midrule\noalign{}
\endhead
\bottomrule\noalign{}
\endlastfoot
Compute & O(n\(^2\) \(\cdot\) d) per layer --- quadratic in sequence length \\
\rowcolor{zebra}
Memory & O(n\(^2\)) for attention matrix --- the KV cache grows with context \\
\end{longtable}

This quadratic bottleneck is why long context is hard and why Flash Attention, linear attention, and sparse attention matter.

\subsection{Multi-Head Attention}\label{multi-head-attention}

Run self-attention h times in parallel, each with different learned projections:

\setcodelang{Python}

\begin{Shaded}
\begin{Highlighting}[]
\NormalTok{head\_i }\OperatorTok{=}\NormalTok{ Attention(Q W\_Q}\OperatorTok{\^{}}\NormalTok{i, K W\_K}\OperatorTok{\^{}}\NormalTok{i, V W\_V}\OperatorTok{\^{}}\NormalTok{i)}

\NormalTok{MultiHead(Q, K, V) }\OperatorTok{=}\NormalTok{ Concat(head\_1, ..., head\_h) W\_O}

\NormalTok{where W\_Q}\OperatorTok{\^{}}\NormalTok{i ∈ ℝ}\OperatorTok{\^{}}\NormalTok{\{d\_model×d\_k\}, d\_k }\OperatorTok{=}\NormalTok{ d\_model }\OperatorTok{/}\NormalTok{ h}
\end{Highlighting}
\end{Shaded}

\textbf{Why multiple heads?}

\begin{itemize}
\tightlist
\item
  Each head can specialize in different relationships
\item
  Head 1 might focus on syntactic dependencies (subject\(\rightarrow\)verb)
\item
  Head 2 might focus on coreference (pronoun\(\rightarrow\)antecedent)
\item
  Head 3 might focus on semantic similarity
\item
  More expressive than a single attention computation
\item
  In practice, different heads empirically learn different linguistic features
\end{itemize}

\textbf{Key numbers}: GPT-3 uses h=96 heads, d\_model=12288, d\_k=128 per head.

\subsection{Positional Encoding}\label{positional-encoding}

Attention is \textbf{permutation-invariant} --- without position information, "dog bites man" = "man bites dog". We inject position information by adding positional encodings to token embeddings.

\subsubsection{Sinusoidal (original Transformer)}\label{sinusoidal-original-transformer}

\setcodelang{}

\begin{verbatim}
PE(pos, 2i)   = sin(pos / 10000^(2i/d_model))
PE(pos, 2i+1) = cos(pos / 10000^(2i/d_model))
\end{verbatim}

\begin{itemize}
\tightlist
\item
  Different frequencies for different dimensions
\item
  Relative positions can be computed from these (key property)
\item
  Generalizes to sequences longer than training length (theoretically)
\end{itemize}

\subsubsection{Learned Absolute Positional Embeddings (GPT-2, BERT)}\label{learned-absolute-positional-embeddings-gpt-2-bert}

Simply learn a vector for each position 0..max\_len. Simple but doesn\textquotesingle t extrapolate.

\subsubsection{Rotary Positional Embedding (RoPE) --- Llama, GPT-NeoX}\label{rotary-positional-embedding-rope--llama-gpt-neox}

Encodes position by rotating Q and K vectors. Relative position naturally falls out of the dot product. Better length generalization.

\subsubsection{ALiBi (PaLM, BLOOM)}\label{alibi-palm-bloom}

Add a linear bias to attention scores based on distance: score -= m * \textbar i - j\textbar. Simple, extrapolates well, no learned parameters.

\subsection{The Full Transformer Block}\label{the-full-transformer-block}

\visualfigure{../figures/aiml-04-tf/02-transformer-block.pdf}{A transformer layer = self-attention (mix information across tokens) + a position-wise MLP, each wrapped in a residual connection and LayerNorm. Stack N of them.}

\subsubsection{Encoder Block}\label{encoder-block}

\setcodelang{}

\begin{verbatim}
Input X
  │
  ├─> LayerNorm → Multi-Head Self-Attention → dropout
  │                                              │
  └─────────────────────────────────────────────+ (residual)
  │
  ├─> LayerNorm → Feed-Forward Network → dropout
  │                                         │
  └─────────────────────────────────────────+ (residual)
  │
Output
\end{verbatim}

\textbf{Feed-Forward Network (FFN):}

\setcodelang{}

\begin{verbatim}
FFN(x) = max(0, x W_1 + b_1) W_2 + b_2

Typical: d_model=512 → d_ff=2048 → d_model=512
         (4x expansion, then project back)
\end{verbatim}

In modern models, often uses GELU or SwiGLU activation instead of ReLU.

\subsubsection{Layer Normalization}\label{layer-normalization}

\setcodelang{}

\begin{verbatim}
LayerNorm(x) = γ * (x - μ) / (σ + ε) + β

where μ, σ computed across features (not batch), γ, β are learned
\end{verbatim}

\textbf{Why not BatchNorm?} Sequences have variable lengths; batch statistics vary with sequence length. LayerNorm normalizes each position independently.

\textbf{Pre-LN vs Post-LN}: original Transformer uses Post-LN (after residual). Modern models use Pre-LN (before sublayer) --- more stable training for deep models.

\subsubsection{Residual Connections}\label{residual-connections}

\setcodelang{Python}

\begin{Shaded}
\begin{Highlighting}[]
\NormalTok{output }\OperatorTok{=}\NormalTok{ LayerNorm(x }\OperatorTok{+}\NormalTok{ Sublayer(x))}
\end{Highlighting}
\end{Shaded}

\begin{itemize}
\tightlist
\item
  Allow gradients to flow directly through the network
\item
  Enable very deep stacking (GPT-3: 96 layers)
\item
  Without residuals, 96-layer Transformers would be untrainable
\end{itemize}

\subsection{Encoder-Only vs Decoder-Only vs Encoder-Decoder}\label{encoder-only-vs-decoder-only-vs-encoder-decoder}

\begin{longtable}[]{@{}
  >{\raggedright\arraybackslash}p{(\columnwidth - 6\tabcolsep) * \real{0.1980}}
  >{\raggedright\arraybackslash}p{(\columnwidth - 6\tabcolsep) * \real{0.2574}}
  >{\raggedright\arraybackslash}p{(\columnwidth - 6\tabcolsep) * \real{0.2376}}
  >{\raggedright\arraybackslash}p{(\columnwidth - 6\tabcolsep) * \real{0.3069}}@{}}
\toprule\noalign{}
\rowcolor{acc!16}
\begin{minipage}[b]{\linewidth}\raggedright
Property
\end{minipage} & \begin{minipage}[b]{\linewidth}\raggedright
Encoder-only (BERT)
\end{minipage} & \begin{minipage}[b]{\linewidth}\raggedright
Decoder-only (GPT)
\end{minipage} & \begin{minipage}[b]{\linewidth}\raggedright
Encoder-Decoder (T5)
\end{minipage} \\
\midrule\noalign{}
\endhead
\bottomrule\noalign{}
\endlastfoot
Attention type & Bidirectional & Causal (left-only) & Cross-attention \\
\rowcolor{zebra}
Training objective & Masked LM (MLM) & Next-token prediction & Seq2seq (span corruption) \\
What it sees & Full context & Only past tokens & Encoder: full; Decoder: causal \\
\rowcolor{zebra}
Best for & Classification, NER, QA & Generation, chat, code & Translation, summarization \\
Examples & BERT, RoBERTa, ELECTRA & GPT-2/3/4, Llama, Claude & T5, BART, mT5 \\
\rowcolor{zebra}
Embedding quality & Excellent & Good & Good \\
Generation & Poor (not designed for it) & Excellent & Good \\
\rowcolor{zebra}
Parameter efficiency & High for understanding & High for generation & Somewhat redundant for pure gen \\
\end{longtable}

\textbf{Why decoder-only won the LLM race?}

\begin{itemize}
\tightlist
\item
  Simpler architecture (one module type)
\item
  Next-token prediction is a universal pretraining objective
\item
  Scales better --- RLHF and instruction tuning work well
\item
  Inference is natural: just keep generating
\end{itemize}

\subsection{Pretraining Objectives}\label{pretraining-objectives}

\subsubsection{Causal Language Modeling (CLM) --- GPT family}\label{causal-language-modeling-clm--gpt-family}

Predict the next token given all previous tokens:

\setcodelang{}

\begin{verbatim}
L_CLM = -Σ_t log P(x_t | x_1, ..., x_{t-1})

"The cat sat on the ___"  →  "mat"
\end{verbatim}

\begin{itemize}
\tightlist
\item
  Self-supervised: targets come from the input itself
\item
  The model learns the full data distribution P(X)
\item
  Generation is natural: sample from P(x\_t \textbar{} x\_\{\textless t\}) autoregressively
\end{itemize}

\subsubsection{Masked Language Modeling (MLM) --- BERT}\label{masked-language-modeling-mlm--bert}

Randomly mask 15\% of tokens, predict the masked tokens:

\setcodelang{}

\begin{verbatim}
"The cat [MASK] on the mat"  →  predict "sat"

L_MLM = -Σ_{masked t} log P(x_t | x_{1..n, t≠masked})
\end{verbatim}

\begin{itemize}
\tightlist
\item
  Bidirectional context \(\rightarrow\) better representations
\item
  Weaker generative model (not trained to generate sequences autoregressively)
\end{itemize}

\subsubsection{Span Corruption (T5)}\label{span-corruption-t5}

Mask contiguous spans of tokens, predict the entire span:

\setcodelang{}

\begin{verbatim}
Input:  "The cat <X> the mat <Y> quickly"
Target: "<X> sat on <Y> moved"
\end{verbatim}

Bridges MLM and CLM --- good for seq2seq tasks.

\subsection{Scaling Laws \& Emergent Abilities}\label{scaling-laws--emergent-abilities}

\subsubsection{Chinchilla Scaling Laws (Hoffmann et al., 2022)}\label{chinchilla-scaling-laws-hoffmann-et-al-2022}

The optimal number of training tokens is \textbf{\textasciitilde20\(\times\) the number of parameters}:

\setcodelang{}

\begin{verbatim}
N_opt ∝ C^0.5  (optimal parameters given compute budget C)
D_opt ∝ C^0.5  (optimal tokens given compute budget C)

Rule of thumb: D_opt ≈ 20 × N_opt
\end{verbatim}

\textbf{Implication}: GPT-3 (175B params, 300B tokens) was undertrained. Chinchilla (70B params, 1.4T tokens) outperformed it. Llama uses this insight.

\subsubsection{Power Law Scaling}\label{power-law-scaling}

\setcodelang{Python}

\begin{Shaded}
\begin{Highlighting}[]
\NormalTok{Loss ∝ (}\DecValTok{1}\OperatorTok{/}\NormalTok{N)}\OperatorTok{\^{}}\NormalTok{α }\OperatorTok{+}\NormalTok{ (}\DecValTok{1}\OperatorTok{/}\NormalTok{D)}\OperatorTok{\^{}}\NormalTok{β }\OperatorTok{+}\NormalTok{ L\_∞}

\NormalTok{where N }\OperatorTok{=}\NormalTok{ model size, D }\OperatorTok{=}\NormalTok{ dataset size, α ≈ β ≈ }\FloatTok{0.5}
\end{Highlighting}
\end{Shaded}

Each decade of compute reduces loss predictably --- no diminishing returns (so far).

\subsubsection{Emergent Abilities}\label{emergent-abilities}

Some capabilities appear suddenly above a threshold model size:

\begin{itemize}
\tightlist
\item
  \textbf{Chain-of-thought reasoning}: \textasciitilde100B params
\item
  \textbf{In-context learning} (few-shot): \textasciitilde10B params
\item
  \textbf{Instruction following}: enabled by RLHF (scale alone not enough)
\item
  \textbf{Multi-step arithmetic}: appears around GPT-3 scale
\end{itemize}

\textbf{Controversy}: Some argue emergence is an artifact of metrics (using non-linear/discontinuous eval metrics). Gradual improvement in underlying probabilities manifests as sudden capability jumps.

\subsection{Context Window \& KV Cache}\label{context-window--kv-cache}

\subsubsection{Context Window}\label{context-window}

The maximum number of tokens the model can attend to simultaneously:

\begin{itemize}
\tightlist
\item
  GPT-3: 2048 tokens
\item
  GPT-4: 8K \(\rightarrow\) 32K \(\rightarrow\) 128K tokens (different versions)
\item
  Claude 3.5: 200K tokens
\item
  Gemini 1.5 Pro: 1M tokens
\end{itemize}

\textbf{Why is long context hard?}

\begin{itemize}
\tightlist
\item
  Attention is O(n\(^2\)) in time and memory
\item
  KV cache grows linearly: storing K, V for all tokens
\item
  Training on long sequences requires massive GPU memory
\item
  Models may not actually use distant context effectively ("lost in the middle" problem)
\end{itemize}

\subsubsection{KV Cache}\label{kv-cache}

During inference (generation), we avoid recomputing K and V for previous tokens:

\setcodelang{}

\begin{verbatim}
Without KV cache (step t):
  Compute Q, K, V for ALL tokens 1..t
  Cost: O(t²) per step, O(n³) total for n tokens

With KV cache (step t):
  Load cached K, V for tokens 1..t-1
  Compute Q, K, V only for new token t
  Cost: O(t) per step, O(n²) total — but memory cost O(n·d·h·L)
\end{verbatim}

\textbf{KV cache memory formula:}

\setcodelang{Python}

\begin{Shaded}
\begin{Highlighting}[]
\NormalTok{Memory }\OperatorTok{=} \DecValTok{2}\NormalTok{ × n\_layers × n\_heads × d\_head × seq\_len × bytes\_per\_element}
       \OperatorTok{=} \DecValTok{2}\NormalTok{ × L × H × d\_k × T × dtype\_bytes}

\NormalTok{For Llama }\DecValTok{2} \DecValTok{70}\ErrorTok{B} \ControlFlowTok{with} \DecValTok{128}\ErrorTok{K}\NormalTok{ context }\KeywordTok{in}\NormalTok{ FP16:}
  \OperatorTok{=} \DecValTok{2}\NormalTok{ × }\DecValTok{80}\NormalTok{ × }\DecValTok{64}\NormalTok{ × }\DecValTok{128}\NormalTok{ × }\DecValTok{131072}\NormalTok{ × }\DecValTok{2} \BuiltInTok{bytes}\NormalTok{ ≈ }\DecValTok{320}\ErrorTok{GB}\NormalTok{ just }\ControlFlowTok{for}\NormalTok{ KV cache}\OperatorTok{!}
\end{Highlighting}
\end{Shaded}

This is why long context is expensive and why techniques like:

\begin{itemize}
\tightlist
\item
  \textbf{Grouped Query Attention (GQA)}: multiple query heads share one K/V pair
\item
  \textbf{Multi-Query Attention (MQA)}: all query heads share a single K/V pair
\item
  \textbf{PagedAttention} (vLLM): page KV cache like OS virtual memory
\end{itemize}

\subsection{Fine-Tuning}\label{fine-tuning}

\subsubsection{Full Fine-Tuning}\label{full-fine-tuning}

Update all parameters on task-specific data:

\begin{itemize}
\tightlist
\item
  Best performance on specific tasks
\item
  Very expensive: same cost as pretraining (GPU-hours, memory)
\item
  Risk of catastrophic forgetting
\item
  Used when you have lots of task data and resources
\end{itemize}

\subsubsection{Instruction Tuning}\label{instruction-tuning}

Fine-tune on (instruction, response) pairs to follow human directions:

\setcodelang{}

\begin{verbatim}
Prompt:  "Summarize this article in 3 bullet points: [article]"
Target:  "• Key point 1\n• Key point 2\n• Key point 3"
\end{verbatim}

\begin{itemize}
\tightlist
\item
  Makes models much more useful without massive compute
\item
  InstructGPT, FLAN-T5, Alpaca used this approach
\item
  Usually followed by RLHF for alignment
\end{itemize}

\subsubsection{PEFT: Parameter-Efficient Fine-Tuning}\label{peft-parameter-efficient-fine-tuning}

Fine-tune only a small fraction of parameters:

\begin{longtable}[]{@{}
  >{\raggedright\arraybackslash}p{(\columnwidth - 4\tabcolsep) * \real{0.2189}}
  >{\raggedright\arraybackslash}p{(\columnwidth - 4\tabcolsep) * \real{0.3098}}
  >{\raggedright\arraybackslash}p{(\columnwidth - 4\tabcolsep) * \real{0.4714}}@{}}
\toprule\noalign{}
\rowcolor{acc!16}
\begin{minipage}[b]{\linewidth}\raggedright
Method
\end{minipage} & \begin{minipage}[b]{\linewidth}\raggedright
Trainable Params
\end{minipage} & \begin{minipage}[b]{\linewidth}\raggedright
How
\end{minipage} \\
\midrule\noalign{}
\endhead
\bottomrule\noalign{}
\endlastfoot
LoRA & \textasciitilde0.1-1\% & Low-rank weight matrices \\
\rowcolor{zebra}
QLoRA & \textasciitilde0.1-1\% & LoRA on quantized base model \\
Adapters & \textasciitilde1-5\% & Small modules between layers \\
\rowcolor{zebra}
Prefix tuning & \textasciitilde0.1\% & Learned prefix tokens \\
Prompt tuning & \textless\textless0.1\% & Soft prompt tokens \\
\end{longtable}

\subsubsection{LoRA (Low-Rank Adaptation)}\label{lora-low-rank-adaptation}

\textbf{Key insight}: weight updates during fine-tuning lie in a low-rank subspace.

Instead of updating W \(\in\) \(\mathbb{R}\)\^{}\{d\(\times\)k\} directly, parameterize the update as:

\setcodelang{Python}

\begin{Shaded}
\begin{Highlighting}[]
\NormalTok{W\_new }\OperatorTok{=}\NormalTok{ W\_pretrained }\OperatorTok{+}\NormalTok{ ΔW }\OperatorTok{=}\NormalTok{ W\_pretrained }\OperatorTok{+}\NormalTok{ B·A}

\NormalTok{where B ∈ ℝ}\OperatorTok{\^{}}\NormalTok{\{d×r\}, A ∈ ℝ}\OperatorTok{\^{}}\NormalTok{\{r×k\}, rank r }\OperatorTok{\textless{}\textless{}} \BuiltInTok{min}\NormalTok{(d, k)}

\NormalTok{Typical: d}\OperatorTok{=}\DecValTok{4096}\NormalTok{, k}\OperatorTok{=}\DecValTok{4096}\NormalTok{, r}\OperatorTok{=}\DecValTok{8} \KeywordTok{or} \DecValTok{16}
\NormalTok{Original params: }\DecValTok{4096}\ErrorTok{²} \OperatorTok{=} \FloatTok{16.7}\ErrorTok{M}
\NormalTok{LoRA params:     }\DecValTok{4096}\NormalTok{×}\DecValTok{8} \OperatorTok{+} \DecValTok{8}\NormalTok{×}\DecValTok{4096} \OperatorTok{=} \DecValTok{65}\ErrorTok{K}\NormalTok{  (}\FloatTok{99.6}\OperatorTok{\%}\NormalTok{ reduction}\OperatorTok{!}\NormalTok{)}
\end{Highlighting}
\end{Shaded}

\textbf{Training:}

\begin{itemize}
\tightlist
\item
  W\_pretrained is frozen
\item
  Only A and B are updated
\item
  At inference: W\_new = W\_pretrained + B\(\cdot\)A (merge or keep separate)
\item
  Initialization: A \textasciitilde{} N(0, \(\sigma\)\(^2\)), B = 0 (so \(\Delta\)W = 0 at start)
\item
  Scale factor \(\alpha\): effective update = (\(\alpha\)/r) \(\cdot\) B\(\cdot\)A
\end{itemize}

\textbf{QLoRA}: Quantize the base model to 4-bit (reducing memory 4-8x), then apply LoRA adapters in 16-bit. Enables fine-tuning 65B models on a single A100 80GB.

\subsection{Alignment: RLHF, Reward Models, PPO, DPO}\label{alignment-rlhf-reward-models-ppo-dpo}

\subsubsection{The Alignment Problem}\label{the-alignment-problem}

A language model pretrained on internet text will:

\begin{itemize}
\tightlist
\item
  Produce harmful content (toxic text exists online)
\item
  Be sycophantic or deceptive
\item
  Not follow instructions consistently
\item
  Optimize for perplexity, not helpfulness
\end{itemize}

We need to align model behavior with human preferences.

\subsubsection{RLHF Pipeline (Reinforcement Learning from Human Feedback)}\label{rlhf-pipeline-reinforcement-learning-from-human-feedback}

Three stages:

\textbf{Stage 1: Supervised Fine-Tuning (SFT)}

\setcodelang{}

\begin{verbatim}
Human writes ideal responses → fine-tune model on these
Result: SFT model — follows instructions but not yet aligned
\end{verbatim}

\textbf{Stage 2: Reward Model (RM) Training}

\setcodelang{}

\begin{verbatim}
- Show human raters K responses to a prompt
- They rank them: response_A > response_B > response_C
- Train RM: r_θ(prompt, response) → scalar reward
- RM loss: -log(σ(r_θ(prompt, chosen) - r_θ(prompt, rejected)))
\end{verbatim}

\textbf{Stage 3: PPO (Proximal Policy Optimization)}

\setcodelang{}

\begin{verbatim}
- Use RM as reward signal to optimize the language model
- Constraint: KL divergence from SFT model (don't drift too far)
- Reward: r_RM(prompt, response) - β * KL(π_RL || π_SFT)
- PPO clip: prevents too-large policy updates
\end{verbatim}

\setcodelang{Python}

\begin{Shaded}
\begin{Highlighting}[]
\NormalTok{PPO objective:}
\NormalTok{L\_PPO }\OperatorTok{=}\NormalTok{ E[}\BuiltInTok{min}\NormalTok{(r\_t(θ)·A\_t, clip(r\_t(θ), }\DecValTok{1}\OperatorTok{{-}}\NormalTok{ε, }\DecValTok{1}\OperatorTok{+}\NormalTok{ε)·A\_t)] }\OperatorTok{{-}}\NormalTok{ β·KL}

\NormalTok{where r\_t(θ) }\OperatorTok{=}\NormalTok{ π\_θ(a}\OperatorTok{|}\NormalTok{s) }\OperatorTok{/}\NormalTok{ π\_θ\_old(a}\OperatorTok{|}\NormalTok{s)  (probability ratio)}
\NormalTok{      A\_t }\OperatorTok{=}\NormalTok{ advantage (reward }\OperatorTok{{-}}\NormalTok{ baseline)}
\end{Highlighting}
\end{Shaded}

\textbf{RLHF challenges:}

\begin{itemize}
\tightlist
\item
  Expensive: requires many human raters
\item
  Reward hacking: model learns to game the RM
\item
  Instability: PPO is notoriously finicky to tune
\item
  Reward model can be wrong
\end{itemize}

\subsubsection{DPO (Direct Preference Optimization)}\label{dpo-direct-preference-optimization}

\textbf{Key insight}: skip the explicit reward model. Directly optimize on preference data.

\setcodelang{Python}

\begin{Shaded}
\begin{Highlighting}[]
\NormalTok{DPO Loss:}
\NormalTok{L\_DPO(θ) }\OperatorTok{=} \OperatorTok{{-}}\NormalTok{E\_\{(x,y\_w,y\_l)\} [}
\NormalTok{    log σ(β · log(π\_θ(y\_w}\OperatorTok{|}\NormalTok{x)}\OperatorTok{/}\NormalTok{π\_ref(y\_w}\OperatorTok{|}\NormalTok{x))}
           \OperatorTok{{-}}\NormalTok{ β · log(π\_θ(y\_l}\OperatorTok{|}\NormalTok{x)}\OperatorTok{/}\NormalTok{π\_ref(y\_l}\OperatorTok{|}\NormalTok{x)))}
\NormalTok{]}

\NormalTok{where y\_w }\OperatorTok{=}\NormalTok{ preferred (winner) response}
\NormalTok{      y\_l }\OperatorTok{=}\NormalTok{ rejected (loser) response}
\NormalTok{      π\_ref }\OperatorTok{=}\NormalTok{ reference policy (SFT model, frozen)}
\NormalTok{      β }\OperatorTok{=}\NormalTok{ temperature controlling KL penalty}
\end{Highlighting}
\end{Shaded}

\textbf{DPO vs RLHF comparison:}

\begin{longtable}[]{@{}
  >{\raggedright\arraybackslash}p{(\columnwidth - 4\tabcolsep) * \real{0.2222}}
  >{\raggedright\arraybackslash}p{(\columnwidth - 4\tabcolsep) * \real{0.4444}}
  >{\raggedright\arraybackslash}p{(\columnwidth - 4\tabcolsep) * \real{0.3333}}@{}}
\toprule\noalign{}
\rowcolor{acc!16}
\begin{minipage}[b]{\linewidth}\raggedright
Aspect
\end{minipage} & \begin{minipage}[b]{\linewidth}\raggedright
RLHF + PPO
\end{minipage} & \begin{minipage}[b]{\linewidth}\raggedright
DPO
\end{minipage} \\
\midrule\noalign{}
\endhead
\bottomrule\noalign{}
\endlastfoot
Complexity & High (3 stages, RM + RL) & Low (1 SFT + 1 stage) \\
\rowcolor{zebra}
Stability & Tricky (PPO hyperparams) & More stable \\
Reward model & Required & Not needed \\
\rowcolor{zebra}
Performance & State of the art & Competitive with RLHF \\
Memory & High (4 models at once) & Lower (2 models) \\
\rowcolor{zebra}
Reward hacking & Possible & Less likely \\
Flexibility & Can adjust reward on the fly & Fixed to dataset \\
\end{longtable}

\textbf{DPO intuition}: implicitly defines a reward model as log ratio of policy to reference. The optimal policy that maximizes this reward under KL constraint has a closed form --- DPO directly optimizes it.

\subsection{Decoding / Inference}\label{decoding--inference}

Given a trained model, how do we generate text?

\subsubsection{Greedy Decoding}\label{greedy-decoding}

\setcodelang{Python}

\begin{Shaded}
\begin{Highlighting}[]
\NormalTok{x\_t }\OperatorTok{=}\NormalTok{ argmax P(x\_t }\OperatorTok{|}\NormalTok{ x\_\{}\OperatorTok{\textless{}}\NormalTok{t\})}
\end{Highlighting}
\end{Shaded}

Always pick the highest-probability token. Fast but repetitive and suboptimal.

\subsubsection{Beam Search}\label{beam-search}

Maintain k (beam width) hypothesis sequences:

\setcodelang{}

\begin{verbatim}
At each step:
  Expand each of k sequences with top-k tokens
  Score all k² candidates by cumulative log-prob
  Keep top k
\end{verbatim}

\begin{itemize}
\tightlist
\item
  Better than greedy for structured outputs (translation)
\item
  Tends to produce generic, boring text for open-ended generation
\item
  Length normalization needed: score / \textbar sequence\textbar\^{}\(\alpha\)
\end{itemize}

\subsubsection{Temperature Sampling}\label{temperature-sampling}

\setcodelang{}

\begin{verbatim}
P_scaled(x) = softmax(logits / T)

T → 0:  approaches greedy (sharp distribution)
T = 1:  original model distribution
T > 1:  more uniform/random (creative but less coherent)
\end{verbatim}

\subsubsection{Top-k Sampling}\label{top-k-sampling}

\setcodelang{Python}

\begin{Shaded}
\begin{Highlighting}[]
\NormalTok{Sample }\ImportTok{from}\NormalTok{ top k tokens only (renormalize the k probabilities)}

\NormalTok{Typical: k }\OperatorTok{=} \DecValTok{40}\OperatorTok{{-}}\DecValTok{200}
\NormalTok{Problem: k may be too large when distribution }\KeywordTok{is}\NormalTok{ flat, too small when peaked}
\end{Highlighting}
\end{Shaded}

\subsubsection{Top-p / Nucleus Sampling (Holtzman et al., 2020)}\label{top-p--nucleus-sampling-holtzman-et-al-2020}

\setcodelang{Python}

\begin{Shaded}
\begin{Highlighting}[]
\NormalTok{Find smallest }\BuiltInTok{set}\NormalTok{ S where Σ\_\{x ∈ S\} P(x) ≥ p}
\NormalTok{Sample }\ImportTok{from}\NormalTok{ S (renormalized)}

\NormalTok{Typical: p }\OperatorTok{=} \FloatTok{0.9} \KeywordTok{or} \FloatTok{0.95}
\NormalTok{Advantage: dynamically adapts to distribution shape}
\end{Highlighting}
\end{Shaded}

\subsubsection{Top-k vs Top-p Comparison}\label{top-k-vs-top-p-comparison}

\begin{longtable}[]{@{}
  >{\raggedright\arraybackslash}p{(\columnwidth - 4\tabcolsep) * \real{0.4504}}
  >{\raggedright\arraybackslash}p{(\columnwidth - 4\tabcolsep) * \real{0.2586}}
  >{\raggedright\arraybackslash}p{(\columnwidth - 4\tabcolsep) * \real{0.2909}}@{}}
\toprule\noalign{}
\rowcolor{acc!16}
\begin{minipage}[b]{\linewidth}\raggedright
Scenario
\end{minipage} & \begin{minipage}[b]{\linewidth}\raggedright
Top-k (k=50)
\end{minipage} & \begin{minipage}[b]{\linewidth}\raggedright
Top-p (p=0.9)
\end{minipage} \\
\midrule\noalign{}
\endhead
\bottomrule\noalign{}
\endlastfoot
Peaked distribution (one clear next word) & Keeps 49 garbage options & Correctly uses just top \textasciitilde5 \\
\rowcolor{zebra}
Flat distribution (many valid continuations) & May be too restrictive & Correctly samples from \textasciitilde200 \\
Recommendation & Less adaptive & More principled \\
\end{longtable}

\subsubsection{Repetition Penalty}\label{repetition-penalty}

\setcodelang{Python}

\begin{Shaded}
\begin{Highlighting}[]
\NormalTok{P\_penalized(x) }\OperatorTok{=}\NormalTok{ P(x) }\OperatorTok{/}\NormalTok{ penalty  }\ControlFlowTok{if}\NormalTok{ x appeared recently}
\NormalTok{                 P(x)             otherwise}

\NormalTok{Typical: penalty }\OperatorTok{=} \FloatTok{1.1}\NormalTok{ to }\FloatTok{1.3}
\end{Highlighting}
\end{Shaded}

\subsubsection{Min-p Sampling (newer)}\label{min-p-sampling-newer}

Only sample from tokens with P(x) \textgreater{} min\_p \(\times\) P(x\_max). Cleaner than top-k.

\subsubsection{Speculative Decoding (inference acceleration)}\label{speculative-decoding-inference-acceleration}

Use a small fast draft model to propose k tokens, then use the large model to verify all k in parallel (one forward pass). Accept tokens where they agree. \textasciitilde3x speedup for free.

\subsection{Quantization}\label{quantization}

Running 70B parameter models requires 140GB in FP16 --- doesn\textquotesingle t fit on commodity hardware. Quantization reduces precision:

\begin{longtable}[]{@{}
  >{\raggedright\arraybackslash}p{(\columnwidth - 8\tabcolsep) * \real{0.1603}}
  >{\raggedright\arraybackslash}p{(\columnwidth - 8\tabcolsep) * \real{0.0819}}
  >{\raggedright\arraybackslash}p{(\columnwidth - 8\tabcolsep) * \real{0.2253}}
  >{\raggedright\arraybackslash}p{(\columnwidth - 8\tabcolsep) * \real{0.2867}}
  >{\raggedright\arraybackslash}p{(\columnwidth - 8\tabcolsep) * \real{0.2458}}@{}}
\toprule\noalign{}
\rowcolor{acc!16}
\begin{minipage}[b]{\linewidth}\raggedright
Format
\end{minipage} & \begin{minipage}[b]{\linewidth}\raggedright
Bits
\end{minipage} & \begin{minipage}[b]{\linewidth}\raggedright
Bytes/param
\end{minipage} & \begin{minipage}[b]{\linewidth}\raggedright
70B model size
\end{minipage} & \begin{minipage}[b]{\linewidth}\raggedright
Quality loss
\end{minipage} \\
\midrule\noalign{}
\endhead
\bottomrule\noalign{}
\endlastfoot
FP32 & 32 & 4B & 280GB & Baseline \\
\rowcolor{zebra}
FP16/BF16 & 16 & 2B & 140GB & Minimal \\
INT8 & 8 & 1B & 70GB & Small \\
\rowcolor{zebra}
INT4 & 4 & 0.5B & 35GB & Noticeable \\
INT2 & 2 & 0.25B & 17.5GB & Significant \\
\end{longtable}

\subsubsection{GPTQ (Post-Training Quantization)}\label{gptq-post-training-quantization}

\begin{itemize}
\tightlist
\item
  Quantize weights to INT4 using second-order information (Hessian)
\item
  Compensate for quantization error layer by layer
\item
  One-shot, no retraining needed
\item
  Used for: running Llama 70B on a single RTX 4090
\end{itemize}

\subsubsection{AWQ (Activation-aware Weight Quantization)}\label{awq-activation-aware-weight-quantization}

\begin{itemize}
\tightlist
\item
  Identifies salient weights (high activation magnitude)
\item
  Protects important weights from quantization
\item
  Better than GPTQ in practice for extreme quantization
\end{itemize}

\subsubsection{GGUF / llama.cpp}\label{gguf--llamacpp}

\begin{itemize}
\tightlist
\item
  Format for CPU inference
\item
  Mixed precision per layer
\item
  Enables running quantized models on MacBook Pro (Apple Silicon)
\end{itemize}

\textbf{Key interview point}: quantization is essential for inference economics. A 4-bit quantized model can run at 25-50\% the cost of FP16 with 5-15\% quality degradation.

\subsection{Mixture of Experts (MoE)}\label{mixture-of-experts-moe}

Instead of using all parameters for every token, use a \textbf{gating network} to route each token to a subset of "expert" feed-forward layers:

\setcodelang{Python}

\begin{Shaded}
\begin{Highlighting}[]
\NormalTok{Standard FFN: y }\OperatorTok{=}\NormalTok{ FFN(x)   [}\BuiltInTok{all}\NormalTok{ params used]}

\NormalTok{MoE FFN:}
\NormalTok{  gates }\OperatorTok{=}\NormalTok{ softmax(W\_gate · x)         }\CommentTok{\# routing weights}
\NormalTok{  top}\OperatorTok{{-}}\NormalTok{K experts selected}
\NormalTok{  y }\OperatorTok{=}\NormalTok{ Σ\_\{i }\KeywordTok{in}\NormalTok{ top}\OperatorTok{{-}}\NormalTok{K\} gates\_i · FFN\_i(x)   }\CommentTok{\# weighted sum}
\end{Highlighting}
\end{Shaded}

\textbf{Key metrics:}

\begin{itemize}
\tightlist
\item
  Total parameters: all experts combined (e.g., 1.8T for GPT-4 estimate)
\item
  Active parameters: experts used per token (e.g., \textasciitilde200B for GPT-4 estimate)
\item
  Efficiency: 8 experts, top-2 \(\rightarrow\) same FLOPs as dense model with 1/4 the params but 4x the capacity
\end{itemize}

\textbf{MoE benefits:}

\begin{itemize}
\tightlist
\item
  Increase model capacity without proportional compute increase
\item
  Experts specialize in different content types
\item
  Same inference cost as smaller dense model
\end{itemize}

\textbf{MoE challenges:}

\begin{itemize}
\tightlist
\item
  Load balancing: all tokens going to same expert \(\rightarrow\) waste
\item
  Training instability: routing collapse possible
\item
  All experts must fit in memory (just not all active simultaneously)
\item
  Communication overhead in distributed settings
\end{itemize}

\textbf{Models using MoE:}

\begin{itemize}
\tightlist
\item
  GPT-4 (estimated \textasciitilde1.8T total, \textasciitilde200B active)
\item
  Mixtral 8x7B (8 experts of 7B each, 2 active \(\rightarrow\) 12.9B active params)
\item
  Google Switch Transformer, GLaM
\item
  Grok-1 (314B total, 86B active)
\end{itemize}

\subsection{Hallucination --- Why It Happens}\label{hallucination--why-it-happens}

\textbf{Hallucination}: the model generates text that is factually incorrect, fabricated, or not grounded in reality, but does so confidently.

\subsubsection{Root causes}\label{root-causes}

\begin{enumerate}
\def\labelenumi{\arabic{enumi}.}
\item
  \textbf{Next-token prediction objective} doesn\textquotesingle t care about truth, only fluency. The model is optimized to produce plausible continuations, not accurate ones.
\item
  \textbf{No grounding mechanism}: knowledge is stored in weights, not retrieved from a verified source. Weights encode patterns, not facts with citations.
\item
  \textbf{Distributional mismatch}: training data has patterns like "The capital of France is {[}most common completion{]}" --- model learns statistical associations, not truth.
\item
  \textbf{Overconfident generation}: temperature sampling and beam search both can confidently produce wrong text when the training distribution doesn\textquotesingle t cover the query.
\item
  \textbf{Long-tail knowledge}: rare facts have weak signal in training data. Model interpolates/extrapolates incorrectly.
\item
  \textbf{RLHF pressure}: reward models may prefer confident-sounding answers \(\rightarrow\) model learns to sound confident even when uncertain.
\end{enumerate}

\subsubsection{Why it\textquotesingle s fundamental (not a bug to fix)}\label{why-its-fundamental-not-a-bug-to-fix}

\begin{itemize}
\tightlist
\item
  The model fundamentally lacks a "truth module" --- it\textquotesingle s a function approximator over text distributions
\item
  Adding retrieval (RAG) helps but doesn\textquotesingle t eliminate it (model can still misinterpret retrieved text)
\item
  Calibration: we want P(model confident) \(\approx\) P(model correct), currently poorly calibrated
\end{itemize}

\subsubsection{Mitigation strategies}\label{mitigation-strategies}

\begin{itemize}
\tightlist
\item
  \textbf{RAG (Retrieval-Augmented Generation)}: ground answers in retrieved documents
\item
  \textbf{Citation}: force model to quote sources
\item
  \textbf{Self-consistency}: sample multiple times, take majority vote
\item
  \textbf{Chain-of-thought}: reasoning steps reduce hallucination on structured tasks
\item
  \textbf{RLHF with fact-checking}: reward truthfulness explicitly
\item
  \textbf{Tool use}: let model call APIs/databases instead of generating facts
\end{itemize}

\section{Architecture / Diagrams}\label{architecture--diagrams}

\subsection{Scaled Dot-Product Attention}\label{scaled-dot-product-attention-1}

\setcodelang{}

\begin{verbatim}
         ┌─────────────────────────────────────────────────────┐
         │                                                     │
  Q ────►│  MatMul(Q, Kᵀ)                                     │
  K ────►│       │                                             │
         │       ▼                                             │
         │  Scale (/√d_k)                                      │
         │       │                                             │
         │       ▼                                             │
         │  [Optional Mask]  ◄── (causal: mask future tokens)  │
         │       │                                             │
         │       ▼                                             │
         │   Softmax                                           │
         │       │                                             │
  V ────►│  MatMul(·, V)                                       │
         │       │                                             │
         │       ▼                                             │
         │   Output                                            │
         └─────────────────────────────────────────────────────┘

Attention(Q,K,V) = softmax( Q·Kᵀ / √d_k ) · V
\end{verbatim}

\subsection{Multi-Head Attention}\label{multi-head-attention-1}

\setcodelang{}

\begin{verbatim}
Input X  ──────────────────────────────────────────────────────────►
         │                    │                    │
         ▼                    ▼                    ▼
   [W_Q1,W_K1,W_V1]   [W_Q2,W_K2,W_V2]   [W_Qh,W_Kh,W_Vh]
         │                    │                    │
         ▼                    ▼                    ▼
   Attention head 1    Attention head 2  ...  Attention head h
         │                    │                    │
         └────────────────────┴────────────────────┘
                              │
                          Concat
                              │
                             W_O
                              │
                           Output
\end{verbatim}

\subsection{Full Transformer Encoder Stack}\label{full-transformer-encoder-stack}

\setcodelang{}

\begin{verbatim}
Input Tokens
     │
     ▼
Token Embeddings + Positional Encodings
     │
     ▼
┌─────────────────────────────────────┐
│         Encoder Block × N           │
│                                     │
│   ┌──────────────────────────────┐  │
│   │   Multi-Head Self-Attention  │  │
│   └──────────────────────────────┘  │
│              │                      │
│          Add & Norm                 │
│              │                      │
│   ┌──────────────────────────────┐  │
│   │    Feed-Forward Network      │  │
│   │  (Linear → GELU → Linear)   │  │
│   └──────────────────────────────┘  │
│              │                      │
│          Add & Norm                 │
└─────────────────────────────────────┘
     │
     ▼
Encoder Output (contextual embeddings for all positions)
\end{verbatim}

\subsection{Full Transformer Decoder Stack (Encoder-Decoder)}\label{full-transformer-decoder-stack-encoder-decoder}

\setcodelang{}

\begin{verbatim}
Target Tokens (shifted right)
     │
     ▼
Token Embeddings + Positional Encodings
     │
     ▼
┌─────────────────────────────────────────────────┐
│              Decoder Block × N                  │
│                                                 │
│  ┌──────────────────────────────────────────┐  │
│  │    Masked Multi-Head Self-Attention       │  │
│  │    (causal: can't see future tokens)      │  │
│  └──────────────────────────────────────────┘  │
│                      │                          │
│                  Add & Norm                     │
│                      │                          │
│  ┌──────────────────────────────────────────┐  │
│  │    Multi-Head Cross-Attention             │  │
│  │    Q from decoder, K/V from encoder       │  │
│  └──────────────────────────────────────────┘  │
│                      │                          │
│                  Add & Norm                     │
│                      │                          │
│  ┌──────────────────────────────────────────┐  │
│  │         Feed-Forward Network             │  │
│  └──────────────────────────────────────────┘  │
│                      │                          │
│                  Add & Norm                     │
└─────────────────────────────────────────────────┘
     │
     ▼
Linear + Softmax → Probability over vocabulary
\end{verbatim}

\subsection{Decoder-Only GPT Block (most modern LLMs)}\label{decoder-only-gpt-block-most-modern-llms}

\setcodelang{}

\begin{verbatim}
Input Tokens
     │
     ▼
Embeddings + Positional Encoding
     │
     ▼
┌─────────────────────────────────────┐
│           GPT Block × L             │
│                                     │
│  ┌──────────────────────────────┐   │
│  │   LayerNorm (Pre-LN)        │   │
│  └──────────────────────────────┘   │
│               │                     │
│  ┌──────────────────────────────┐   │
│  │ Causal Multi-Head Attention  │   │
│  │ (mask: each token sees only  │   │
│  │  tokens to its left)         │   │
│  └──────────────────────────────┘   │
│               │                     │
│          + Residual                 │
│               │                     │
│  ┌──────────────────────────────┐   │
│  │   LayerNorm (Pre-LN)        │   │
│  └──────────────────────────────┘   │
│               │                     │
│  ┌──────────────────────────────┐   │
│  │    Feed-Forward Network      │   │
│  │    (SwiGLU / GELU)          │   │
│  └──────────────────────────────┘   │
│               │                     │
│          + Residual                 │
└─────────────────────────────────────┘
     │
     ▼
LayerNorm (final)
     │
     ▼
Linear (d_model → vocab_size)
     │
     ▼
Softmax → Sample next token
     │
     ▼
Append token, repeat →
\end{verbatim}

\subsection{Positional Encoding Visualization}\label{positional-encoding-visualization}

\setcodelang{}

\begin{verbatim}
Position  Dim 0    Dim 1    Dim 2    Dim 3
  0       sin(0)   cos(0)   sin(0)   cos(0)    = [0,  1,  0,   1  ]
  1       sin(1)   cos(1)   sin(.1)  cos(.1)   = [.84,.54,.10, .99]
  2       sin(2)   cos(2)   sin(.2)  cos(.2)   = [.91,-.41,.20,.98]
  ...

Low-frequency dims capture large-scale position structure
High-frequency dims capture fine-grained position
\end{verbatim}

\subsection{RLHF Pipeline}\label{rlhf-pipeline}

\setcodelang{}

\begin{verbatim}
STAGE 1: Supervised Fine-Tuning (SFT)
─────────────────────────────────────
Human expert writes ideal responses
         │
         ▼
   Fine-tune base LM ──► SFT Model (π_SFT)


STAGE 2: Reward Model Training
──────────────────────────────
Prompt ──► SFT Model ──► 4 responses
                               │
                     Human ranking (A>B>C>D)
                               │
                    Train Reward Model (r_θ)
                    Loss: log σ(r(y_w) - r(y_l))


STAGE 3: PPO Fine-Tuning
────────────────────────
Prompt ──► RL Policy (π_θ) ──► Response
                                    │
                         Reward Model (r_θ) ──► Score
                                    │
                  KL(π_θ || π_SFT) ──► KL Penalty
                                    │
                         PPO Update ──► Better π_θ
                         (keep θ close to SFT to avoid reward hacking)


RESULT: Helpful, Harmless, Honest model (HHH)
\end{verbatim}

\subsection{Tokenization Example}\label{tokenization-example}

\setcodelang{}

\begin{verbatim}
Input text:  "ChatGPT is transformative!"

BPE Tokenization (GPT-4 tiktoken):
  "Chat"     → ID: 14126
  "G"        → ID: 38
  "PT"       → ID: 2898
  " is"      → ID: 374
  " transform" → ID: 5276
  "ative"    → ID: 1413
  "!"        → ID: 0

Token count: 7 (vs 4 words)
Note: space before "is" included in token " is"
Note: punctuation is separate token
\end{verbatim}

\section{Real-World Examples}\label{real-world-examples}

\subsection{ChatGPT / GPT-4 (OpenAI)}\label{chatgpt--gpt-4-openai}

\begin{itemize}
\tightlist
\item
  \textbf{Architecture}: decoder-only Transformer (estimated \textasciitilde1.8T params MoE for GPT-4)
\item
  \textbf{Training}: internet text \(\rightarrow\) SFT \(\rightarrow\) RLHF
\item
  \textbf{Use cases}: Q\&A, coding, writing, reasoning
\item
  \textbf{Context}: GPT-4o supports 128K tokens
\item
  \textbf{Impact}: 100M users in 60 days, \$1B+ ARR
\end{itemize}

\subsection{GitHub Copilot (Microsoft/OpenAI)}\label{github-copilot-microsoftopenai}

\begin{itemize}
\tightlist
\item
  Powered by Codex (GPT-4 class model fine-tuned on code)
\item
  Context: your current file + open tabs + recently viewed files
\item
  Generates whole functions, docstrings, tests
\item
  \$19/month, millions of developers \(\rightarrow\) \textasciitilde30\% of Copilot-written code accepted
\end{itemize}

\subsection{Retrieval-Augmented Generation (RAG)}\label{retrieval-augmented-generation-rag}

\setcodelang{}

\begin{verbatim}
User query ──► Embedding model ──► Vector search over document store
                                         │
                               Top-k documents retrieved
                                         │
                    [Document 1][Document 2][Document 3][Query]
                                         │
                               LLM generates answer
                               grounded in retrieved docs
\end{verbatim}

Used in: Bing Chat, Perplexity, enterprise Q\&A systems

\subsection{Google Search / AI Overviews}\label{google-search--ai-overviews}

\begin{itemize}
\tightlist
\item
  BERT used for query understanding since 2019 (biggest leap in 5 years)
\item
  Gemini powers AI Overviews (generative answers at top of results)
\item
  8.5B searches/day --- even small improvements matter enormously
\end{itemize}

\subsection{Claude (Anthropic)}\label{claude-anthropic}

\begin{itemize}
\tightlist
\item
  Constitutional AI: instead of human rating every pair, give the model a "constitution" of principles and have it rate its own outputs
\item
  Long context specialist: 200K tokens --- can read entire codebases
\item
  Computer Use: multimodal model that can control a computer
\end{itemize}

\subsection{Llama (Meta AI)}\label{llama-meta-ai}

\begin{itemize}
\tightlist
\item
  Open weights: anyone can run/fine-tune/distill
\item
  Enabled explosion of open-source ecosystem: Vicuna, Alpaca, Mistral, etc.
\item
  Llama 3.1 405B competes with GPT-4 on many benchmarks
\item
  Powers Meta AI on WhatsApp, Instagram, Facebook
\end{itemize}

\section{Real-Life Analogies --- The Busy Newsroom}\label{real-life-analogies--the-busy-newsroom}

\emph{Imagine a newsroom where editors collaboratively read and rewrite a breaking story. Every element of Transformer LLMs maps naturally to this setting.}

\begin{longtable}[]{@{}
  >{\raggedright\arraybackslash}p{(\columnwidth - 2\tabcolsep) * \real{0.2098}}
  >{\raggedright\arraybackslash}p{(\columnwidth - 2\tabcolsep) * \real{0.7902}}@{}}
\toprule\noalign{}
\rowcolor{acc!16}
\begin{minipage}[b]{\linewidth}\raggedright
Concept
\end{minipage} & \begin{minipage}[b]{\linewidth}\raggedright
Newsroom Analogy
\end{minipage} \\
\midrule\noalign{}
\endhead
\bottomrule\noalign{}
\endlastfoot
\textbf{Tokens} & Individual words written on sticky notes across the story board \\
\rowcolor{zebra}
\textbf{Self-attention} & Each editor draws colored string from their word to every other word they find relevant \\
\textbf{Query (Q)} & The question an editor writes on their pad: "Who is the subject here?" \\
\rowcolor{zebra}
\textbf{Key (K)} & The label tag each sticky note carries: "I am a person\textquotesingle s name", "I am a verb" \\
\textbf{Value (V)} & The actual content written on the sticky note that gets passed along \\
\rowcolor{zebra}
\textbf{Attention weights} & How thick/bright each string is --- more relevant words get stronger connections \\
\textbf{Multi-head attention} & Several editors working the board simultaneously, each focused on a different angle: one tracks grammar, one tracks factual consistency, one tracks tone \\
\rowcolor{zebra}
\textbf{Positional encoding} & The page number and line number stamped on each sticky note --- without it, "Dog bites man" = "Man bites dog" \\
\textbf{Residual connections} & Each editor\textquotesingle s draft preserves the previous draft underneath --- changes are additive, not destructive \\
\rowcolor{zebra}
\textbf{Layer normalization} & The copy desk\textquotesingle s style guide that standardizes formatting after each pass \\
\textbf{Transformer layers} & Successive editing passes --- draft 1, draft 2, draft 3 --- each round refines and adds meaning \\
\rowcolor{zebra}
\textbf{Feed-forward network} & Each editor privately re-processes their word\textquotesingle s meaning before the next round of attention \\
\textbf{Pretraining} & The editor spent 20 years reading every newspaper, book, and website ever published \\
\rowcolor{zebra}
\textbf{Fine-tuning} & Two weeks of house-style training for \emph{this} publication --- learn the voice, the format, the audience \\
\textbf{RLHF} & An editor-in-chief sits with the team, rates their drafts, and gives detailed feedback --- the team adjusts \\
\rowcolor{zebra}
\textbf{DPO} & Instead of the EIC rating drafts, they provide "prefer A over B" comparisons; the team infers the preference directly \\
\textbf{Temperature} & How boldly an editor paraphrases --- low temp: safe, literal; high temp: creative, risky \\
\rowcolor{zebra}
\textbf{Top-p sampling} & The editor only considers rewordings that, together, cover 90\% of plausible options \\
\textbf{Context window} & How many sticky notes fit on the editor\textquotesingle s desk at once --- the rest are in a filing cabinet they can\textquotesingle t see \\
\rowcolor{zebra}
\textbf{KV cache} & The editor\textquotesingle s running notes on previously read paragraphs --- no need to re-read them \\
\textbf{Hallucination} & An editor confidently inventing a quote from a source they never actually read \\
\rowcolor{zebra}
\textbf{MoE (Mixture of Experts)} & A pool of 100 specialist editors; the routing desk sends each sticky note to the 2 most relevant specialists \\
\textbf{LoRA} & Teaching an editor just 1\% new habits for this client --- personality stays intact, style adjusts \\
\rowcolor{zebra}
\textbf{Quantization} & Compressing detailed editorial notes to shorthand symbols --- faster to read, slight accuracy loss \\
\textbf{RAG} & Before writing, the editor runs to the library and pulls the three most relevant reference articles \\
\rowcolor{zebra}
\textbf{Beam search} & Keeping the top 5 draft directions simultaneously, then picking the best final story \\
\textbf{Tokenizer} & The editor\textquotesingle s typesetting team that breaks words into standard-sized type blocks \\
\end{longtable}

\section{Memory Tricks / Mnemonics}\label{memory-tricks--mnemonics}

\subsection{QKV --- "Ask, Label, Content"}\label{qkv--ask-label-content}

\begin{itemize}
\tightlist
\item
  \textbf{Q}uery = \textbf{Ask}: what is this token looking for?
\item
  \textbf{K}ey = \textbf{Label}: what is this token advertising?
\item
  \textbf{V}alue = \textbf{Content}: what does this token actually contribute?
\end{itemize}

\subsection{"Scale to Prevent Flat Softmax"}\label{scale-to-prevent-flat-softmax}

The scaling factor \(\surd\)d\_k prevents attention scores from growing too large \(\rightarrow\) softmax saturation \(\rightarrow\) vanishing gradients. Remember: \textbf{large scores \(\rightarrow\) flat gradients \(\rightarrow\) stuck training}.

\subsection{BERT vs GPT --- "Understand vs Generate"}\label{bert-vs-gpt--understand-vs-generate}

\begin{itemize}
\tightlist
\item
  \textbf{B}ERT = \textbf{B}idirectional = \textbf{B}etter for understanding (can see left AND right)
\item
  \textbf{G}PT = \textbf{G}enerator = always looks \textbf{l}eft (causal masking)
\end{itemize}

\subsection{LoRA = "Frozen Elephant, Tiny Tail"}\label{lora--frozen-elephant-tiny-tail}

The pretrained model (elephant) is frozen. You only train a tiny adapter (tail). The elephant still runs the circus; you just teach it a few new tricks.

\subsection{\texorpdfstring{RLHF stages --- "SFT \(\rightarrow\) RM \(\rightarrow\) PPO"}{RLHF stages --- "SFT \textbackslash rightarrow RM \textbackslash rightarrow PPO"}}\label{rlhf-stages--sft-rightarrow-rm-rightarrow-ppo}

\textbf{S}oft \textbf{R}obot \textbf{P}retty

\begin{itemize}
\tightlist
\item
  \textbf{S}FT: teach the robot to follow instructions
\item
  \textbf{R}M: teach it what "pretty" means (via human ratings)
\item
  \textbf{P}PO: reinforce the pretty behaviors
\end{itemize}

\subsection{\texorpdfstring{Scaling laws --- "Compute = Parameters \(\times\) Tokens"}{Scaling laws --- "Compute = Parameters \textbackslash times Tokens"}}\label{scaling-laws--compute--parameters-times-tokens}

Chinchilla says: \textbf{N \(\times\) 20 = D} (optimal tokens \(\approx\) 20\(\times\) params). If you have a 10B model, train on 200B tokens. If you trained on less, you left performance on the table.

\subsection{\texorpdfstring{Hallucination --- "Fluent \(\neq\) Factual"}{Hallucination --- "Fluent \textbackslash neq Factual"}}\label{hallucination--fluent-neq-factual}

The model is trained to predict the next token, not to be correct. \textbf{It\textquotesingle s a distribution over text, not a database of facts.}

\subsection{Attention complexity --- "N-squared nightmare"}\label{attention-complexity--n-squared-nightmare}

Attention is O(n\(^2\)) in sequence length. Double the context \(\rightarrow\) quadruple the compute. This is why GPT-3 used 2K context and getting to 1M context required years of research.

\subsection{Temperature extremes}\label{temperature-extremes}

\begin{itemize}
\tightlist
\item
  \textbf{T \(\rightarrow\) 0} = greedy (max probability always)
\item
  \textbf{T = 1} = raw model (original distribution)
\item
  \textbf{T \textgreater{} 1} = wild/creative (more uniform)
\item
  \textbf{T = 0.7} = typical creative use; \textbf{T = 0.0} = factual/deterministic
\end{itemize}

\subsection{\texorpdfstring{Decoding methods --- "Greedy \(\rightarrow\) Beam \(\rightarrow\) Sample"}{Decoding methods --- "Greedy \textbackslash rightarrow Beam \textbackslash rightarrow Sample"}}\label{decoding-methods--greedy-rightarrow-beam-rightarrow-sample}

More creativity from left to right:

\setcodelang{}

\begin{verbatim}
Greedy < Beam Search < Top-k < Top-p (nucleus) < High Temperature
[deterministic]                                    [creative/random]
\end{verbatim}

\section{Common Interview Questions}\label{common-interview-questions}

\subsection{Q1: Explain self-attention from first principles}\label{q1-explain-self-attention-from-first-principles}

\textbf{Model Answer:}

Self-attention allows each token in a sequence to compute a weighted sum of all other tokens\textquotesingle{} representations, where weights reflect semantic relevance.

For each token, we compute three vectors via learned linear projections: Query (what am I looking for?), Key (what do I offer?), Value (what is my content?).

The attention scores are computed as: \texttt{A\ =\ softmax(Q·Kᵀ\ /\ √d\_k)\ ·\ V}

The \texttt{Q·Kᵀ} gives an n\(\times\)n matrix of raw compatibilities. We divide by \(\surd\)d\_k to prevent softmax saturation --- without scaling, large d\_k causes dot products to grow large, pushing gradients toward zero. Softmax normalizes each row to get attention weights (probabilities). The final output is a weighted sum of Values.

\textbf{Key properties}: permutation equivariant (order-agnostic without positional encodings), global receptive field (every token attends to every other), differentiable and trainable end-to-end.

\textbf{Follow-up}: Why is it called "self" attention? \(\rightarrow\) Q, K, V all come from the SAME input sequence (vs cross-attention where Q comes from the decoder and K/V from the encoder).

\subsection{\texorpdfstring{Q2: Why do we scale by \(\surd\)d\_k?}{Q2: Why do we scale by \textbackslash surdd\_k?}}\label{q2-why-do-we-scale-by-surdd_k}

\textbf{Model Answer:}

Consider Q and K vectors with i.i.d. components from N(0,1). The dot product q\(\cdot\)k = \(\Sigma\)\(_i\) q\(_i\)k\(_i\) has expected value 0 and \textbf{variance d\_k} (sum of d\_k independent unit-variance products). So std(q\(\cdot\)k) = \(\surd\)d\_k.

When d\_k is large (e.g., 128), raw dot products have std \textasciitilde11. After softmax, most probability mass concentrates on a single token \(\rightarrow\) gradients vanish (softmax in saturation region).

Dividing by \(\surd\)d\_k normalizes dot products to std\textasciitilde1, keeping softmax in a reasonable gradient-friendly regime.

\textbf{Interviewers love this because} it shows you understand the math, not just the formula.

\subsection{Q3: What is RLHF and why is it needed?}\label{q3-what-is-rlhf-and-why-is-it-needed}

\textbf{Model Answer:}

RLHF (Reinforcement Learning from Human Feedback) is a three-stage process to align language models with human preferences:

\begin{enumerate}
\def\labelenumi{\arabic{enumi}.}
\tightlist
\item
  \textbf{SFT}: Fine-tune the pretrained model on human-written ideal responses
\item
  \textbf{Reward Model}: Train a model to predict human preference scores from comparison data (human says response A \textgreater{} response B \(\rightarrow\) train RM to give A higher score)
\item
  \textbf{PPO}: Use the RM as a reward signal to further fine-tune the language model via RL. Add KL divergence penalty from SFT model to prevent reward hacking.
\end{enumerate}

\textbf{Why needed?} A model trained purely on next-token prediction optimizes for \emph{fluency and perplexity}, not helpfulness or safety. It will happily generate toxic content, refuse to admit uncertainty, or provide dangerous information if that\textquotesingle s what appeared in training data.

RLHF bridges the gap between "good at predicting text" and "good at being a helpful assistant."

\textbf{Follow-up}: What\textquotesingle s the main failure mode? \(\rightarrow\) \textbf{Reward hacking}: the policy learns to generate responses that fool the reward model without being actually better. Solution: KL penalty, better RM, iterative RLHF.

\subsection{Q4: RLHF vs DPO --- when would you choose each?}\label{q4-rlhf-vs-dpo--when-would-you-choose-each}

\textbf{Model Answer:}

\begin{longtable}[]{@{}
  >{\raggedright\arraybackslash}p{(\columnwidth - 4\tabcolsep) * \real{0.1528}}
  >{\raggedright\arraybackslash}p{(\columnwidth - 4\tabcolsep) * \real{0.4444}}
  >{\raggedright\arraybackslash}p{(\columnwidth - 4\tabcolsep) * \real{0.4028}}@{}}
\toprule\noalign{}
\rowcolor{acc!16}
\begin{minipage}[b]{\linewidth}\raggedright
Criterion
\end{minipage} & \begin{minipage}[b]{\linewidth}\raggedright
Choose RLHF
\end{minipage} & \begin{minipage}[b]{\linewidth}\raggedright
Choose DPO
\end{minipage} \\
\midrule\noalign{}
\endhead
\bottomrule\noalign{}
\endlastfoot
Resources & You have large infrastructure & Resource-constrained \\
\rowcolor{zebra}
Flexibility & Need to update reward on the fly & Fixed preference dataset \\
Stability & Can tune PPO carefully & Want stable, simpler training \\
\rowcolor{zebra}
Scale & Large team, production system & Research, smaller teams \\
Performance & Potentially higher ceiling & Competitive in practice \\
\end{longtable}

\textbf{DPO insight}: it can be shown that under certain assumptions, the optimal policy satisfying a KL-penalized reward objective has a closed form. DPO directly optimizes this --- the log ratio of policy to reference implicitly represents the reward. No separate RM needed.

\textbf{When DPO fails}: when the preference data doesn\textquotesingle t cover the distribution the policy explores. RLHF with online data collection can handle this; DPO is an offline algorithm.

\subsection{Q5: What is LoRA and why does it work?}\label{q5-what-is-lora-and-why-does-it-work}

\textbf{Model Answer:}

LoRA (Low-Rank Adaptation) fine-tunes LLMs by adding low-rank decompositions to weight matrices instead of updating weights directly.

For weight matrix W \(\in\) \(\mathbb{R}\)\^{}\{d\(\times\)k\}, the update \(\Delta\)W is parameterized as:

\setcodelang{Python}

\begin{Shaded}
\begin{Highlighting}[]
\NormalTok{ΔW }\OperatorTok{=}\NormalTok{ B · A    where B ∈ ℝ}\OperatorTok{\^{}}\NormalTok{\{d×r\}, A ∈ ℝ}\OperatorTok{\^{}}\NormalTok{\{r×k\}, r }\OperatorTok{\textless{}\textless{}} \BuiltInTok{min}\NormalTok{(d,k)}
\end{Highlighting}
\end{Shaded}

W\_pretrained is \textbf{frozen}; only A and B are trained.

\textbf{Why it works --- the hypothesis}: fine-tuning updates lie in a low-rank subspace. The "intrinsic dimensionality" of adaptation is small --- you don\textquotesingle t need to update all O(d\(^2\)) parameters to specialize a model.

\textbf{Practical benefits}:

\begin{itemize}
\tightlist
\item
  1000\(\times\) fewer trainable parameters (r=8 for a 4096\(\times\)4096 layer = 65K vs 16.7M)
\item
  Can train 7B model on a single consumer GPU
\item
  Adapters can be stored and swapped cheaply
\item
  Can merge: W\_merged = W\_pretrained + B\(\cdot\)A (zero inference overhead)
\end{itemize}

\textbf{QLoRA extension}: quantize W\_pretrained to 4-bit (saving 4\(\times\) memory), keep A and B in 16-bit. Enables fine-tuning 65B models on a single A100.

\subsection{Q6: What is the KV cache and why does it matter?}\label{q6-what-is-the-kv-cache-and-why-does-it-matter}

\textbf{Model Answer:}

During autoregressive inference, the model generates tokens one at a time. At step t, it needs attention over all t tokens. Without caching, we recompute K and V for tokens 1..t-1 at every step \(\rightarrow\) O(n\(^3\)) total complexity.

The KV cache stores the K and V vectors for all past tokens. At each new step, we only compute Q, K, V for the new token, retrieve cached K and V, and compute attention \(\rightarrow\) O(n\(^2\)) total complexity.

\textbf{Cost}: KV cache memory = \texttt{2\ ×\ L\ ×\ H\ ×\ d\_k\ ×\ T\ ×\ bytes\_per\_element}

For Llama 2 70B serving with 128K context in FP16:

\begin{itemize}
\tightlist
\item
  2 \(\times\) 80 layers \(\times\) 64 heads \(\times\) 128 d\_head \(\times\) 131072 tokens \(\times\) 2 bytes \(\approx\) 320GB
\end{itemize}

This often exceeds GPU memory per request. Solutions:

\begin{itemize}
\tightlist
\item
  \textbf{GQA/MQA}: share K/V across query heads
\item
  \textbf{PagedAttention} (vLLM): virtual memory for KV cache, allows memory sharing across requests
\item
  \textbf{Quantized KV cache}: store K/V in INT8
\end{itemize}

\textbf{Interview point}: KV cache is a key driver of inference cost and a critical consideration in production serving.

\subsection{Q7: What are scaling laws? What\textquotesingle s the Chinchilla finding?}\label{q7-what-are-scaling-laws-whats-the-chinchilla-finding}

\textbf{Model Answer:}

Scaling laws (Kaplan et al., 2020 originally) show that LLM performance scales as a power law with model size N, dataset size D, and compute C:

\setcodelang{Python}

\begin{Shaded}
\begin{Highlighting}[]
\NormalTok{Loss ∝ N}\OperatorTok{\^{}}\NormalTok{\{}\OperatorTok{{-}}\NormalTok{α\} }\OperatorTok{+}\NormalTok{ D}\OperatorTok{\^{}}\NormalTok{\{}\OperatorTok{{-}}\NormalTok{β\} }\OperatorTok{+}\NormalTok{ L\_∞}
\end{Highlighting}
\end{Shaded}

The \textbf{Chinchilla finding} (Hoffmann et al., 2022): given a fixed compute budget, you should equally scale N and D. The optimal ratio is:

\setcodelang{}

\begin{verbatim}
D_opt ≈ 20 × N_opt
\end{verbatim}

This means GPT-3 (175B params, 300B tokens) was \textbf{model-overfit} --- undertrained relative to its size. Chinchilla (70B params, 1.4T tokens) outperformed it with the same compute by using more data.

\textbf{Implications}:

\begin{itemize}
\tightlist
\item
  Llama models are trained following Chinchilla-optimal schedules
\item
  More recent models train for even more tokens (Llama 3: 15T tokens on 8B model \(\rightarrow\) 1875\(\times\) params)
\item
  Training beyond Chinchilla optimal can still improve if you care about inference cost (pay once for training, many times for inference)
\end{itemize}

\section{Senior-Level Discussion Points}\label{senior-level-discussion-points}

\subsection{1. KV Cache Economics at Scale}\label{1-kv-cache-economics-at-scale}

\setcodelang{}

\begin{verbatim}
Per-request cost breakdown (GPT-4 scale estimate):
  - Prefill (processing input): O(n × d × L) FLOPs — compute-bound
  - Decode (generating output): O(1 × d × L + KV cache load) per token — memory-bound

At long contexts (128K):
  - KV cache dominates memory bandwidth
  - 320GB+ per user for 70B model at 128K context
  - Multi-tenancy requires PagedAttention (vLLM) to share GPU memory
\end{verbatim}

\textbf{Business impact}: context length costs money. A 128K context request may cost 100\(\times\) a 1K request. This is why per-token pricing matters more than per-request pricing.

\subsection{2. Context Length Scaling --- What\textquotesingle s Hard}\label{2-context-length-scaling--whats-hard}

\begin{itemize}
\tightlist
\item
  \textbf{Attention is O(n\(^2\))}: Flash Attention reduces constant factors but doesn\textquotesingle t change complexity
\item
  \textbf{Position generalization}: models trained on 4K context don\textquotesingle t work on 32K without tricks (RoPE scaling, YaRN)
\item
  \textbf{"Lost in the middle"} (Liu et al., 2023): models actually use beginning and end of context better than the middle --- attention to middle tokens is empirically weaker
\item
  \textbf{Memory}: Flash Attention solves GPU memory for training (recompute instead of store attention matrices) but KV cache still needed for inference
\end{itemize}

\subsection{3. MoE Routing Decisions}\label{3-moe-routing-decisions}

Key design choices in MoE:

\begin{itemize}
\tightlist
\item
  \textbf{Top-k routing}: each token routed to k of E experts. Standard is k=2, E=8
\item
  \textbf{Load balancing}: auxiliary loss to prevent all tokens going to same experts
\item
  \textbf{Expert capacity}: each expert has a "capacity factor" --- tokens exceeding capacity are dropped or sent to a backup
\item
  \textbf{Token routing vs expert routing}: MegaBlocks uses expert-parallel approach
\end{itemize}

\textbf{Why MoE is interesting for future scaling}: you can increase total parameters (and thus capacity) without proportionally increasing FLOPs per token. The cost of MoE is communication overhead in distributed settings and memory to hold all experts.

\subsection{4. Why Hallucination Is Fundamental (Not a Bug)}\label{4-why-hallucination-is-fundamental-not-a-bug}

The deeper point: \textbf{language models are not knowledge bases}. They are functions that map contexts to probability distributions over next tokens. Factuality requires:

\begin{enumerate}
\def\labelenumi{\arabic{enumi}.}
\tightlist
\item
  Accurate representation of facts in weights (dependent on training data quality/coverage)
\item
  Retrieval of the right representation at inference time (dependent on attention and activation)
\item
  Calibrated uncertainty (model knowing what it doesn\textquotesingle t know)
\end{enumerate}

None of these are explicitly optimized for. The model learns to \emph{sound} authoritative because authoritative-sounding text has high probability in training data.

RAG helps but doesn\textquotesingle t fully solve it: the model can still misread retrieved text, blend multiple sources incorrectly, or confidently interpolate between retrieved facts.

The fundamental solution may require a paradigm shift: \textbf{verification systems, formal knowledge bases, or models explicitly trained to express uncertainty} rather than pure generative modeling.

\subsection{5. Inference Economics --- Practical Constraints}\label{5-inference-economics--practical-constraints}

\setcodelang{Python}

\begin{Shaded}
\begin{Highlighting}[]
\NormalTok{Batch size considerations:}
  \OperatorTok{{-}}\NormalTok{ Prefill phase: large batch }\OperatorTok{=}\NormalTok{ high GPU utilization (compute}\OperatorTok{{-}}\NormalTok{bound)}
  \OperatorTok{{-}}\NormalTok{ Decode phase: batch size limited by KV cache memory (memory}\OperatorTok{{-}}\NormalTok{bound)}
  \OperatorTok{{-}}\NormalTok{ Optimal: prefill }\KeywordTok{in}\NormalTok{ large batches, decode }\ControlFlowTok{with}\NormalTok{ continuous batching}

\NormalTok{Throughput vs latency tradeoff:}
  \OperatorTok{{-}}\NormalTok{ High throughput: maximize tokens}\OperatorTok{/}\NormalTok{second (batch many requests)}
  \OperatorTok{{-}}\NormalTok{ Low latency: minimize time}\OperatorTok{{-}}\NormalTok{to}\OperatorTok{{-}}\NormalTok{first}\OperatorTok{{-}}\NormalTok{token }\OperatorTok{+}\NormalTok{ per}\OperatorTok{{-}}\NormalTok{token latency}
  \OperatorTok{{-}}\NormalTok{ Speculative decoding can improve both }\ControlFlowTok{for}\NormalTok{ the right use cases}

\NormalTok{Cost structure (approximate }\DecValTok{2024}\NormalTok{ prices):}
  \OperatorTok{{-}}\NormalTok{ A100 }\DecValTok{80}\ErrorTok{GB}\NormalTok{: }\OperatorTok{\textasciitilde{}}\NormalTok{$}\DecValTok{2}\OperatorTok{{-}}\DecValTok{3}\OperatorTok{/}\NormalTok{hour on cloud}
  \OperatorTok{{-}} \DecValTok{70}\ErrorTok{B}\NormalTok{ model }\KeywordTok{in}\NormalTok{ FP16: needs }\DecValTok{2}\NormalTok{×A100 }\ControlFlowTok{for}\NormalTok{ inference}
  \OperatorTok{{-}}\NormalTok{ At }\DecValTok{1000}\NormalTok{ tokens}\OperatorTok{/}\NormalTok{s throughput: }\OperatorTok{\textasciitilde{}}\NormalTok{$}\FloatTok{0.002}\OperatorTok{{-}}\FloatTok{0.003}\NormalTok{ per }\DecValTok{1}\ErrorTok{K}\NormalTok{ tokens}
  \OperatorTok{{-}}\NormalTok{ GPT}\OperatorTok{{-}}\DecValTok{4}\ErrorTok{o}\NormalTok{ pricing: $}\DecValTok{5}\OperatorTok{/}\NormalTok{M }\BuiltInTok{input}\NormalTok{, $}\DecValTok{15}\OperatorTok{/}\NormalTok{M output}
\end{Highlighting}
\end{Shaded}

\section{Typical Mistakes Candidates Make}\label{typical-mistakes-candidates-make}

\subsection{1. Confusing attention types}\label{1-confusing-attention-types}

\begin{itemize}
\tightlist
\item
  "BERT uses causal attention" --- WRONG. BERT uses bidirectional (full) attention with masking for MLM.
\item
  "GPT uses cross-attention between encoder and decoder" --- WRONG. GPT is decoder-only, no encoder, no cross-attention.
\end{itemize}

\subsection{2. Getting the scaling formula wrong}\label{2-getting-the-scaling-formula-wrong}

\begin{itemize}
\tightlist
\item
  Saying "bigger model always better" --- WRONG. Chinchilla: you need 20\(\times\) more tokens than parameters. An undertrained big model loses to a well-trained small model.
\end{itemize}

\subsection{3. Misunderstanding LoRA rank}\label{3-misunderstanding-lora-rank}

\begin{itemize}
\tightlist
\item
  Saying "LoRA uses a low-rank matrix to replace weights" --- WRONG. LoRA \emph{adds} a low-rank update to frozen weights. The base weights are NOT replaced.
\item
  Forgetting B is initialized to zero so \(\Delta\)W=0 at the start.
\end{itemize}

\subsection{4. Confusing temperature and top-p}\label{4-confusing-temperature-and-top-p}

\begin{itemize}
\tightlist
\item
  "Higher temperature means fewer options" --- WRONG. Higher temperature flattens the distribution, MORE tokens become viable.
\item
  Not knowing that temperature and top-p are often used together.
\end{itemize}

\subsection{\texorpdfstring{5. The "why \(\surd\)d\_k" blank}{5. The "why \textbackslash surdd\_k" blank}}\label{5-the-why-surdd_k-blank}

This is a litmus test question. Many candidates know the formula but can\textquotesingle t explain the reasoning. Know: variance of dot product grows with d\_k \(\rightarrow\) softmax saturation \(\rightarrow\) gradient vanishing.

\subsection{6. Misexplaining RLHF as "supervised learning on human-written text"}\label{6-misexplaining-rlhf-as-supervised-learning-on-human-written-text}

RLHF Stage 1 (SFT) IS supervised learning. But Stage 3 (PPO) is RL --- the model generates outputs, gets scored, and updates via policy gradient. Candidates conflate SFT and RLHF.

\subsection{7. Not knowing what KV cache means for memory}\label{7-not-knowing-what-kv-cache-means-for-memory}

Saying "long context is fine, just increase context length" without knowing the memory cost. At 128K tokens, KV cache can be 320GB+ for a 70B model.

\subsection{8. Claiming "hallucination is a data quality problem"}\label{8-claiming-hallucination-is-a-data-quality-problem}

It\textquotesingle s partly that, but more fundamentally: the training objective (next-token prediction) doesn\textquotesingle t optimize for factuality. Even a model trained on perfect data will hallucinate because it generates statistically plausible completions, not verified facts.

\subsection{9. Forgetting the residual connections}\label{9-forgetting-the-residual-connections}

When asked to describe the Transformer block, many candidates say "attention \(\rightarrow\) FFN" and omit Add \& Norm. Residuals are critical for training deep networks.

\subsection{10. Treating BERT and GPT as interchangeable}\label{10-treating-bert-and-gpt-as-interchangeable}

They have fundamentally different architectures, objectives, and use cases. Knowing when to use each (and that most modern LLMs use decoder-only) matters.

\section{How This Connects To Other Topics}\label{how-this-connects-to-other-topics}

\subsection{Deep Learning (03-deep-learning.md)}\label{deep-learning-03-deep-learningmd}

\begin{itemize}
\tightlist
\item
  Transformers use all core DL concepts: backprop, Adam optimizer, batch normalization \(\rightarrow\) Layer Normalization, residual networks
\item
  Attention mechanism is built on learned linear layers (Q/K/V projections are just matrix multiplications)
\item
  GPT training: cross-entropy loss, gradient clipping, learning rate schedules (warmup + cosine decay)
\item
  Mixture of Experts extends the MLP/FFN concept
\end{itemize}

\subsection{LLM Applications / RAG}\label{llm-applications--rag}

\begin{itemize}
\tightlist
\item
  Context window limitations motivate RAG: retrieve relevant documents to fit within context
\item
  Embeddings from encoder models (e.g., text-embedding-ada-002) used for semantic search
\item
  Vector databases (Pinecone, Weaviate, FAISS) store embeddings for retrieval
\item
  Chunking strategies matter: how to split documents for embedding
\end{itemize}

\subsection{MLOps (production deployment)}\label{mlops-production-deployment}

\begin{itemize}
\tightlist
\item
  Model serving: vLLM, TGI (Text Generation Inference), TensorRT-LLM
\item
  Quantization for cost reduction: INT8/INT4 inference
\item
  Autoscaling challenges: GPU provisioning is expensive and slow
\item
  Monitoring: latency (TTFT, TPS), token cost, quality drift
\item
  A/B testing: hard because outputs are open-ended text
\end{itemize}

\subsection{System Design}\label{system-design}

\begin{itemize}
\tightlist
\item
  LLM serving system: load balancer \(\rightarrow\) prefill cluster \(\rightarrow\) decode cluster
\item
  Batch processing vs streaming responses (SSE/WebSockets for real-time)
\item
  Caching: exact-match cache for repeated prompts, semantic cache for similar queries
\item
  Rate limiting: tokens/minute, not just requests/minute
\item
  Cost optimization: smaller models for simple tasks, big models for complex
\end{itemize}

\section{FAANG / AI-Lab Interview Tips}\label{faang--ai-lab-interview-tips}

\subsection{Know the levels}\label{know-the-levels}

\begin{itemize}
\tightlist
\item
  \textbf{L4/E4 (junior ML)}: know attention formula, BERT vs GPT, basic fine-tuning concepts
\item
  \textbf{L5/E5 (mid-level ML)}: deep understanding of attention complexity, RLHF pipeline, LoRA mechanics, inference optimization
\item
  \textbf{L6/E6 (senior ML)}: MoE design, KV cache economics, scaling laws, alignment research tradeoffs, system design for LLM serving
\item
  \textbf{Research roles}: mechanistic interpretability, novel architecture proposals, theoretical understanding
\end{itemize}

\subsection{What top labs actually ask}\label{what-top-labs-actually-ask}

\textbf{OpenAI}: Self-attention derivation from scratch (whiteboard coding), RLHF reward hacking mitigations, how you\textquotesingle d scale pretraining to 10T parameters

\textbf{Anthropic}: Constitutional AI principles, DPO vs RLHF tradeoffs, mechanistic interpretability (circuits), safety considerations in deployment, 200K context serving

\textbf{Google DeepMind}: Chinchilla scaling laws, Flash Attention algorithm, TPU-specific parallelism strategies (tensor parallelism), distributed training (ZeRO stages)

\textbf{Meta AI}: Open-source tradeoffs, Llama architecture decisions (GQA, SwiGLU, RoPE), quantization for on-device models (Llama on phone)

\textbf{Microsoft}: LoRA for enterprise fine-tuning, cost optimization for Azure serving, integration with Office 365

\subsection{Coding questions you should be able to implement}\label{coding-questions-you-should-be-able-to-implement}

\setcodelang{Python}

\begin{Shaded}
\begin{Highlighting}[]
\CommentTok{\# 1. Self{-}attention from scratch}
\KeywordTok{def}\NormalTok{ attention(Q, K, V, d\_k, mask}\OperatorTok{=}\VariableTok{None}\NormalTok{):}
\NormalTok{    scores }\OperatorTok{=}\NormalTok{ Q }\OperatorTok{@}\NormalTok{ K.T }\OperatorTok{/}\NormalTok{ (d\_k }\OperatorTok{**} \FloatTok{0.5}\NormalTok{)}
    \ControlFlowTok{if}\NormalTok{ mask }\KeywordTok{is} \KeywordTok{not} \VariableTok{None}\NormalTok{:}
\NormalTok{        scores }\OperatorTok{=}\NormalTok{ scores.masked\_fill(mask }\OperatorTok{==} \DecValTok{0}\NormalTok{, }\BuiltInTok{float}\NormalTok{(}\StringTok{\textquotesingle{}{-}inf\textquotesingle{}}\NormalTok{))}
\NormalTok{    weights }\OperatorTok{=}\NormalTok{ F.softmax(scores, dim}\OperatorTok{={-}}\DecValTok{1}\NormalTok{)}
    \ControlFlowTok{return}\NormalTok{ weights }\OperatorTok{@}\NormalTok{ V}

\CommentTok{\# 2. Multi{-}head attention (structure)}
\KeywordTok{class}\NormalTok{ MultiHeadAttention(nn.Module):}
    \KeywordTok{def} \FunctionTok{\_\_init\_\_}\NormalTok{(}\VariableTok{self}\NormalTok{, d\_model, n\_heads):}
        \VariableTok{self}\NormalTok{.h }\OperatorTok{=}\NormalTok{ n\_heads}
        \VariableTok{self}\NormalTok{.d\_k }\OperatorTok{=}\NormalTok{ d\_model }\OperatorTok{//}\NormalTok{ n\_heads}
        \VariableTok{self}\NormalTok{.W\_Q }\OperatorTok{=}\NormalTok{ nn.Linear(d\_model, d\_model)}
        \VariableTok{self}\NormalTok{.W\_K }\OperatorTok{=}\NormalTok{ nn.Linear(d\_model, d\_model)}
        \VariableTok{self}\NormalTok{.W\_V }\OperatorTok{=}\NormalTok{ nn.Linear(d\_model, d\_model)}
        \VariableTok{self}\NormalTok{.W\_O }\OperatorTok{=}\NormalTok{ nn.Linear(d\_model, d\_model)}

\CommentTok{\# 3. LoRA layer}
\KeywordTok{class}\NormalTok{ LoRALinear(nn.Module):}
    \KeywordTok{def} \FunctionTok{\_\_init\_\_}\NormalTok{(}\VariableTok{self}\NormalTok{, W\_pretrained, rank}\OperatorTok{=}\DecValTok{8}\NormalTok{, alpha}\OperatorTok{=}\DecValTok{16}\NormalTok{):}
        \VariableTok{self}\NormalTok{.W }\OperatorTok{=}\NormalTok{ W\_pretrained  }\CommentTok{\# frozen}
        \VariableTok{self}\NormalTok{.A }\OperatorTok{=}\NormalTok{ nn.Parameter(torch.randn(rank, W.shape[}\DecValTok{1}\NormalTok{]) }\OperatorTok{*} \FloatTok{0.01}\NormalTok{)}
        \VariableTok{self}\NormalTok{.B }\OperatorTok{=}\NormalTok{ nn.Parameter(torch.zeros(W.shape[}\DecValTok{0}\NormalTok{], rank))}
        \VariableTok{self}\NormalTok{.scale }\OperatorTok{=}\NormalTok{ alpha }\OperatorTok{/}\NormalTok{ rank}
    
    \KeywordTok{def}\NormalTok{ forward(}\VariableTok{self}\NormalTok{, x):}
        \ControlFlowTok{return}\NormalTok{ x }\OperatorTok{@} \VariableTok{self}\NormalTok{.W.T }\OperatorTok{+}\NormalTok{ (x }\OperatorTok{@} \VariableTok{self}\NormalTok{.A.T }\OperatorTok{@} \VariableTok{self}\NormalTok{.B.T) }\OperatorTok{*} \VariableTok{self}\NormalTok{.scale}
\end{Highlighting}
\end{Shaded}

\subsection{Communication tips}\label{communication-tips}

\begin{enumerate}
\def\labelenumi{\arabic{enumi}.}
\tightlist
\item
  \textbf{Draw the diagram first}: for attention and Transformer block questions, start with ASCII or whiteboard drawing
\item
  \textbf{Give the formula, then explain it}: \texttt{Attention\ =\ softmax(QKᵀ/√d\_k)V} \(\rightarrow\) "Let me break down each part..."
\item
  \textbf{Know the "why"}: interviewers at top labs care less about memorization, more about whether you understand the reasoning
\item
  \textbf{Connect to trade-offs}: always mention the cost/benefit (e.g., "multi-head attention increases expressiveness at the cost of the projection overhead")
\item
  \textbf{Honest about uncertainty}: better to say "I believe it\textquotesingle s X but let me think through the math" than to confidently state something wrong
\end{enumerate}

\section{Revision Cheat Sheet}\label{revision-cheat-sheet}

\subsection{10-Minute Summary}\label{10-minute-summary}

\begin{enumerate}
\def\labelenumi{\arabic{enumi}.}
\tightlist
\item
  \textbf{Transformers} replaced RNNs in 2017 by using attention (not recurrence) for sequence processing, enabling parallelization
\item
  \textbf{Self-attention}: for each token, compute weighted sum of all tokens\textquotesingle{} values; weights from Q\(\cdot\)K\(^{\mathsf{T}}\)/\(\surd\)d\_k \(\rightarrow\) softmax
\item
  \textbf{Multi-head}: run h attention heads in parallel, each learning different relationship types
\item
  \textbf{Positional encodings}: sine/cosine or learned vectors added to embeddings to inject position information
\item
  \textbf{Block}: LayerNorm \(\rightarrow\) MultiHead Attention \(\rightarrow\) residual \(\rightarrow\) LayerNorm \(\rightarrow\) FFN \(\rightarrow\) residual
\item
  \textbf{Encoder-only} (BERT): bidirectional, MLM objective, good for understanding tasks
\item
  \textbf{Decoder-only} (GPT): causal masking, next-token prediction, good for generation --- the dominant modern choice
\item
  \textbf{Scaling laws} (Chinchilla): optimal tokens \(\approx\) 20\(\times\) params; bigger isn\textquotesingle t always better if undertrained
\item
  \textbf{RLHF}: SFT \(\rightarrow\) reward model \(\rightarrow\) PPO; aligns model with human preferences
\item
  \textbf{DPO}: directly optimize on preference pairs without explicit RM; simpler, competitive
\item
  \textbf{LoRA}: freeze base weights, add low-rank updates B\(\cdot\)A (r=8-64); 100-1000\(\times\) fewer trainable params
\item
  \textbf{KV cache}: cache K,V for past tokens during inference; crucial for latency but huge memory cost
\item
  \textbf{Hallucination}: fundamental --- model optimizes fluency, not truth; RAG and citations help
\item
  \textbf{MoE}: route tokens to subset of expert FFN layers; scale capacity without scaling FLOPs
\item
  \textbf{Quantization}: INT8/INT4 reduces model size and inference cost with acceptable quality loss
\end{enumerate}

\subsection{Key Formulas}\label{key-formulas}

\setcodelang{Python}

\begin{Shaded}
\begin{Highlighting}[]
\NormalTok{Self}\OperatorTok{{-}}\NormalTok{Attention:     Attn(Q,K,V) }\OperatorTok{=}\NormalTok{ softmax(QKᵀ}\OperatorTok{/}\NormalTok{√d\_k) · V}

\NormalTok{Multi}\OperatorTok{{-}}\NormalTok{Head:         MH(Q,K,V) }\OperatorTok{=}\NormalTok{ Concat(head\_1,...,head\_h) · W\_O}
\NormalTok{                    head\_i }\OperatorTok{=}\NormalTok{ Attn(QW\_Qi, KW\_Ki, VW\_Vi)}

\NormalTok{Positional:         PE(pos,}\DecValTok{2}\ErrorTok{i}\NormalTok{)   }\OperatorTok{=}\NormalTok{ sin(pos}\OperatorTok{/}\DecValTok{10000}\OperatorTok{\^{}}\NormalTok{\{}\DecValTok{2}\ErrorTok{i}\OperatorTok{/}\NormalTok{d\})}
\NormalTok{                    PE(pos,}\DecValTok{2}\ErrorTok{i}\OperatorTok{+}\DecValTok{1}\NormalTok{) }\OperatorTok{=}\NormalTok{ cos(pos}\OperatorTok{/}\DecValTok{10000}\OperatorTok{\^{}}\NormalTok{\{}\DecValTok{2}\ErrorTok{i}\OperatorTok{/}\NormalTok{d\})}

\NormalTok{Chinchilla:         D\_opt ≈ }\DecValTok{20}\NormalTok{ × N\_opt}

\NormalTok{LoRA:               ΔW }\OperatorTok{=}\NormalTok{ B·A,  B∈ℝ}\OperatorTok{\^{}}\NormalTok{\{d×r\}, A∈ℝ}\OperatorTok{\^{}}\NormalTok{\{r×k\}, r}\OperatorTok{\textless{}\textless{}}\BuiltInTok{min}\NormalTok{(d,k)}

\NormalTok{DPO loss:           }\OperatorTok{{-}}\NormalTok{E[log σ(β·log(π(y\_w)}\OperatorTok{/}\NormalTok{π\_ref(y\_w)) }\OperatorTok{{-}}\NormalTok{ β·log(π(y\_l)}\OperatorTok{/}\NormalTok{π\_ref(y\_l)))]}

\NormalTok{Temperature:        P\_T(x) }\OperatorTok{=}\NormalTok{ softmax(logits}\OperatorTok{/}\NormalTok{T)}

\NormalTok{KV cache mem:       }\DecValTok{2}\NormalTok{ × L × H × d\_k × T × dtype\_bytes}
\end{Highlighting}
\end{Shaded}

\subsection{Architecture Comparison Cheat Sheet}\label{architecture-comparison-cheat-sheet}

\begin{longtable}[]{@{}
  >{\raggedright\arraybackslash}p{(\columnwidth - 8\tabcolsep) * \real{0.1940}}
  >{\raggedright\arraybackslash}p{(\columnwidth - 8\tabcolsep) * \real{0.2090}}
  >{\raggedright\arraybackslash}p{(\columnwidth - 8\tabcolsep) * \real{0.1791}}
  >{\raggedright\arraybackslash}p{(\columnwidth - 8\tabcolsep) * \real{0.1791}}
  >{\raggedright\arraybackslash}p{(\columnwidth - 8\tabcolsep) * \real{0.2388}}@{}}
\toprule\noalign{}
\rowcolor{acc!16}
\begin{minipage}[b]{\linewidth}\raggedright
\end{minipage} & \begin{minipage}[b]{\linewidth}\raggedright
BERT
\end{minipage} & \begin{minipage}[b]{\linewidth}\raggedright
GPT-2/3
\end{minipage} & \begin{minipage}[b]{\linewidth}\raggedright
T5
\end{minipage} & \begin{minipage}[b]{\linewidth}\raggedright
Llama 3
\end{minipage} \\
\midrule\noalign{}
\endhead
\bottomrule\noalign{}
\endlastfoot
Type & Encoder-only & Decoder-only & Enc-Dec & Decoder-only \\
\rowcolor{zebra}
Attention & Bidirectional & Causal & Cross-attn & Causal + GQA \\
Pretraining & MLM & CLM & Span corrupt & CLM \\
\rowcolor{zebra}
Positional & Learned abs. & Learned abs. & Relative & RoPE \\
Best for & NLU, embedding & Generation & Seq2seq & Generation, chat \\
\rowcolor{zebra}
Normalization & Post-LN & Post-LN & Pre-LN & RMSNorm Pre-LN \\
\end{longtable}

\subsection{Fine-Tuning Comparison}\label{fine-tuning-comparison}

\begin{longtable}[]{@{}
  >{\raggedright\arraybackslash}p{(\columnwidth - 8\tabcolsep) * \real{0.1795}}
  >{\raggedright\arraybackslash}p{(\columnwidth - 8\tabcolsep) * \real{0.1419}}
  >{\raggedright\arraybackslash}p{(\columnwidth - 8\tabcolsep) * \real{0.1290}}
  >{\raggedright\arraybackslash}p{(\columnwidth - 8\tabcolsep) * \real{0.2131}}
  >{\raggedright\arraybackslash}p{(\columnwidth - 8\tabcolsep) * \real{0.3365}}@{}}
\toprule\noalign{}
\rowcolor{acc!16}
\begin{minipage}[b]{\linewidth}\raggedright
Method
\end{minipage} & \begin{minipage}[b]{\linewidth}\raggedright
Trainable \%
\end{minipage} & \begin{minipage}[b]{\linewidth}\raggedright
GPU Memory
\end{minipage} & \begin{minipage}[b]{\linewidth}\raggedright
Performance
\end{minipage} & \begin{minipage}[b]{\linewidth}\raggedright
Use case
\end{minipage} \\
\midrule\noalign{}
\endhead
\bottomrule\noalign{}
\endlastfoot
Full fine-tuning & 100\% & Very high & Best & Large resources, specific task \\
\rowcolor{zebra}
LoRA & 0.1-1\% & Low & Near-full & Standard fine-tuning \\
QLoRA & 0.1-1\% & Very low & Slightly below LoRA & Consumer GPU fine-tuning \\
\rowcolor{zebra}
Adapters & 1-5\% & Low-medium & Good & Modular task switching \\
Prompt tuning & \textless0.1\% & Minimal & Lower & Very limited resources \\
\end{longtable}

\subsection{Decoding Method Comparison}\label{decoding-method-comparison}

\begin{longtable}[]{@{}
  >{\raggedright\arraybackslash}p{(\columnwidth - 8\tabcolsep) * \real{0.2308}}
  >{\raggedright\arraybackslash}p{(\columnwidth - 8\tabcolsep) * \real{0.1718}}
  >{\raggedright\arraybackslash}p{(\columnwidth - 8\tabcolsep) * \real{0.2580}}
  >{\raggedright\arraybackslash}p{(\columnwidth - 8\tabcolsep) * \real{0.0950}}
  >{\raggedright\arraybackslash}p{(\columnwidth - 8\tabcolsep) * \real{0.2444}}@{}}
\toprule\noalign{}
\rowcolor{acc!16}
\begin{minipage}[b]{\linewidth}\raggedright
Method
\end{minipage} & \begin{minipage}[b]{\linewidth}\raggedright
Determinism
\end{minipage} & \begin{minipage}[b]{\linewidth}\raggedright
Quality
\end{minipage} & \begin{minipage}[b]{\linewidth}\raggedright
Speed
\end{minipage} & \begin{minipage}[b]{\linewidth}\raggedright
Use case
\end{minipage} \\
\midrule\noalign{}
\endhead
\bottomrule\noalign{}
\endlastfoot
Greedy & Fully & Low (repetitive) & Fastest & Simple extraction \\
\rowcolor{zebra}
Beam search (k=5) & Yes & High for structured & Slower & Translation, code \\
Top-k (k=50) & No & Good & Fast & Chat, creative \\
\rowcolor{zebra}
Top-p (p=0.9) & No & Good, adaptive & Fast & General generation \\
Top-p + temp=0.7 & No & Best creative & Fast & Creative writing \\
\end{longtable}

\subsection{RLHF vs DPO Comparison}\label{rlhf-vs-dpo-comparison}

\begin{longtable}[]{@{}
  >{\raggedright\arraybackslash}p{(\columnwidth - 4\tabcolsep) * \real{0.2083}}
  >{\raggedright\arraybackslash}p{(\columnwidth - 4\tabcolsep) * \real{0.4722}}
  >{\raggedright\arraybackslash}p{(\columnwidth - 4\tabcolsep) * \real{0.3194}}@{}}
\toprule\noalign{}
\rowcolor{acc!16}
\begin{minipage}[b]{\linewidth}\raggedright
\end{minipage} & \begin{minipage}[b]{\linewidth}\raggedright
RLHF (PPO)
\end{minipage} & \begin{minipage}[b]{\linewidth}\raggedright
DPO
\end{minipage} \\
\midrule\noalign{}
\endhead
\bottomrule\noalign{}
\endlastfoot
Stages & SFT \(\rightarrow\) RM \(\rightarrow\) PPO & SFT \(\rightarrow\) DPO \\
\rowcolor{zebra}
Reward model & Explicit, trained & Implicit \\
Stability & Low (PPO finicky) & High \\
\rowcolor{zebra}
Complexity & High & Low \\
Online/Offline & Online (samples during training) & Offline (fixed dataset) \\
\rowcolor{zebra}
Memory (models) & 4 (ref, RM, value, policy) & 2 (ref, policy) \\
Best when & Production, large team & Research, smaller teams \\
\end{longtable}

\subsection{Most Important Concepts (ranked by interview frequency)}\label{most-important-concepts-ranked-by-interview-frequency}

\begin{enumerate}
\def\labelenumi{\arabic{enumi}.}
\tightlist
\item
  \textbf{Self-attention mechanism} (Q/K/V, formula, why \(\surd\)d\_k) --- asked at nearly every AI interview
\item
  \textbf{BERT vs GPT architecture difference} --- fundamental knowledge check
\item
  \textbf{RLHF pipeline} (3 stages: SFT \(\rightarrow\) RM \(\rightarrow\) PPO) --- all major AI labs use this
\item
  \textbf{LoRA} (what it is, why it works, the math) --- standard at any fine-tuning question
\item
  \textbf{Scaling laws / Chinchilla} (20\(\times\) rule, compute-optimal) --- shows you understand LLM development
\item
  \textbf{KV cache} (what, why, cost) --- critical for systems/production roles
\item
  \textbf{Hallucination} (root cause, not just symptoms) --- distinguishes junior from senior
\item
  \textbf{DPO vs RLHF} --- hot research topic, shows currency
\item
  \textbf{Tokenization} (BPE, why subwords, implications) --- often overlooked but important
\item
  \textbf{MoE} (routing, active vs total params) --- important for GPT-4-level systems
\end{enumerate}

\emph{Last updated: 2026-06-14 \textbar{} Path: AI-ML/04-transformers-and-llms.md}

\end{document}
